%% Question bank chapter file - generated by bank_tools.py build (do not edit while the bank is still being built:
%% it is regenerated from the per-test data in Question Bank\data). Every question is one bankq block:
%%   \begin{bankq}[id=..., source={...}, paper=.., qno=.., type=..., subtopic={...}, answer={...}, answer_src=key|solved,
%%                 working={...}, difficulty=1-10, diff_parts={K S I C F}, diff_note={...},
%%                 group=<passage id or empty>, also_in={duplicate copies in other tests}]
%% difficulty = min(10, K+S+I+C+F-1): K concepts 1-3, S steps 1-4, I insight 0-2, C computation 0-2, F format 0-1
%%   ...question LaTeX, options, figures (\bfig / \bfigi -> img/)...
%%   \end{bankq}
%% Passages shared by several questions are bankpassage blocks; their questions carry group=<passage id>.
\documentclass[10pt]{article}
\usepackage[a4paper,margin=1.8cm]{geometry}
\usepackage[T1]{fontenc}
\usepackage{amsmath,amssymb,mathtools}
\usepackage{graphicx}
\usepackage[export]{adjustbox}
\usepackage{xcolor,needspace,keyval}
\usepackage{enumitem}
\usepackage{booktabs,array,tabularx}
\usepackage[hidelinks]{hyperref}
\setlength{\parindent}{0pt}
\setlength{\parskip}{3pt}
\makeatletter
\define@key{bq}{id}{\def\bq@id{#1}}\define@key{bq}{source}{\def\bq@src{#1}}\define@key{bq}{paper}{\def\bq@paper{#1}}
\define@key{bq}{qno}{\def\bq@qno{#1}}\define@key{bq}{type}{\def\bq@type{#1}}\define@key{bq}{subtopic}{\def\bq@sub{#1}}
\define@key{bq}{answer}{\def\bq@ans{#1}}\define@key{bq}{answer_src}{\def\bq@asrc{#1}}\define@key{bq}{working}{\def\bq@work{#1}}
\define@key{bq}{group}{\def\bq@group{#1}}\define@key{bq}{also_in}{\def\bq@also{#1}}
\define@key{bq}{difficulty}{\def\bq@diff{#1}}\define@key{bq}{diff_parts}{\def\bq@dparts{#1}}\define@key{bq}{diff_note}{\def\bq@dnote{#1}}
\newenvironment{bankq}[1][]{\def\bq@work{}\def\bq@also{}\def\bq@group{}\def\bq@diff{?}\def\bq@dparts{}\def\bq@dnote{}\setkeys{bq}{#1}%
  \par\needspace{4\baselineskip}\medskip\noindent{\small\textbf{\bq@id}\hfill\color{gray}[\bq@src, Paper \bq@paper, Q\bq@qno]
  \bq@type\ $\cdot$ \bq@sub\ $\cdot$ \textbf{D\bq@diff}}\par\nopagebreak\relax}
  {\par{\small\color{teal}Answer: \textbf{\bq@ans}\ifx\bq@work\@empty\else\ (solved: \bq@work)\fi
  \ifx\bq@also\@empty\else\par\color{gray}Also in: \bq@also\fi}\par}
\define@key{bp}{id}{\def\bp@id{#1}}\define@key{bp}{source}{\def\bp@src{#1}}\define@key{bp}{questions}{\def\bp@qs{#1}}
\newenvironment{bankpassage}[1][]{\setkeys{bp}{#1}\par\needspace{4\baselineskip}\bigskip\noindent
  {\small\textbf{Passage \bp@id}\hfill\color{gray}[\bp@src; questions \bp@qs]}\par\nopagebreak\relax\itshape}{\par}
\makeatother
\newcommand{\bfig}[2]{\par\smallskip{\centering\includegraphics[width=#1pt,max width=0.62\linewidth,max height=0.32\textheight]{img/#2}\par}\smallskip}
\newcommand{\bfigi}[2]{\includegraphics[width=#1pt,max width=\linewidth,max height=0.32\textheight,valign=c]{img/#2}}
\newcommand{\bmissing}[1]{\textit{[figure missing in the source conversion: \detokenize{#1}]}}
\newenvironment{qoptions}{\begin{itemize}[leftmargin=1.8em,itemsep=1pt,topsep=2pt,label={}]}{\end{itemize}}
\newcommand{\opt}[2]{\item[(#1)] #2}
\providecommand{\fwstart}{}\providecommand{\fwend}{}\providecommand{\qfullsplit}{}
\newcolumntype{L}[1]{>{\raggedright\arraybackslash}p{#1}}
\newcolumntype{C}[1]{>{\centering\arraybackslash}p{#1}}
\newcolumntype{R}[1]{>{\raggedleft\arraybackslash}p{#1}}
\providecommand{\sectionhdr}[1]{}\providecommand{\subsectionhdr}[1]{}


\begin{document}
{\Large\bfseries Chemical Equilibrium --- Question Bank}\par\smallskip
{\small Chemistry chapter C06. Questions from all converted test papers, grouped by source series. Answers: printed key where it matches the paper, otherwise solved (answer\_src=solved, with working). Exact repeats across papers appear once, with also\_in listing the other copies.}\par\medskip\hrule

\section*{Score-I papers}
\begin{bankq}[id=SI03-P1-CHEM-S3-Q02, source={Score-I - Test 03 (07 Jan 2022)}, paper=1, qno=2, type=numerical, subtopic={Kc and Kp}, answer={40.00 to 40.20}, answer_src=key, difficulty=7, diff_parts={K2 S2 I2 C1 F1}, diff_note={concepts: equilibrium constant, phosgene decomposition from Gibbs energy | steps: delta G = delta H - T delta S; K = exp(-delta G / RT); evaluate | I: combine enthalpy and entropy data into K}]
Equilibrium constant ($K$) at 27$^\circ$C for the decomposition of poisonous gas COCl$_2$, represented by

COCl$_2$(g) $\rightleftharpoons$ CO(g) $+$ Cl$_2$(g) is $\dfrac{1}{2}\times10^{-7}$ bar

Find magnitude of $\Delta_r G$ (kJ/mole) for the following reaction at 27$^\circ$C

CO(g) $+$ Cl$_2$(g) $\rightleftharpoons$ COCl$_2$(g)

if mixture of 1 mole COCl$_2$(g), 1 mole CO(g) and 2 mole Cl$_2$(g) is taken in closed container and maintained at 1 bar pressure and 27$^\circ$C. ($R=8.3$ J/mole/K)
\end{bankq}
\begin{bankq}[id=SI06-P1-CHEM-S1-Q03, source={Score-I - Test 06 (28 Jan 2022)}, paper=1, qno=3, type=single, subtopic={Kc and Kp}, answer={B}, answer_src=key, difficulty=6, diff_parts={K2 S2 I2 C1 F0}, diff_note={concepts: reaction quotient or direction from Kp, partial pressures | steps: Kp of 43 per atm; reaction quotient from concentrations at the given temperature; compare; relation between Kc and Kp | I: units of Kp and conversion to Kc}]
For a gaseous chemical reaction the value of K$_\mathrm{P}$ = $\dfrac43$atm$^{-1}$ at 227$^\circ$C in a closed container. The correct statement for the reaction at same temperature is/are : (Given R = $\dfrac{1}{12}\dfrac{\text{litre}-\text{atm}}{\text{k}-\text{mol}}$)
\begin{qoptions}
\opt{A}{K$_\mathrm{C}$ = $\dfrac{20}{3}$ mole litre$^{-1}$}
\opt{B}{If volume of container is increased equilibrium shifts in backward direction.}
\opt{C}{If partial pressure of each reactant is doubled K$_\mathrm{P}$ becomes 4 times.}
\opt{D}{Can't measure K$_\mathrm{C}$ for the reaction as data are insufficient.}
\end{qoptions}
\end{bankq}
\begin{bankq}[id=SI06-P1-CHEM-S2-Q11, source={Score-I - Test 06 (28 Jan 2022)}, paper=1, qno=11, type=multi, subtopic={Kc and Kp}, answer={A,C,D}, answer_src=key, difficulty=6, diff_parts={K2 S2 I1 C1 F1}, diff_note={concepts: two simultaneous equilibria, a common product | steps: Kp expressions for PCl5 and COCl2 decompositions sharing chlorine; solve partial pressures; check statements}]
Following two equilibria are simultaneously established in a container

$\mathrm{PCl_5(g)\rightleftharpoons PCl_3(g)+Cl_2}$ ; $\mathrm{CO(g)+Cl_2(g)\rightleftharpoons COCl_2(g)}$

If some Ni(s) is introduced in the container establishing another equilibrium.

Which among the following statements is/are definitely true about this system ?
\begin{qoptions}
\opt{A}{COCl$_2$ concentration decreases as the new equilibrium is established.}
\opt{B}{As Ni is a solid it will not have any impact on the existing two equilibrium reactions.}
\opt{C}{PCl$_3$ concentration decreases to re--attain the equilibrium.}
\opt{D}{PCl$_5$ concentration increases to attain the new equilibrium.}
\end{qoptions}
\end{bankq}
\begin{bankq}[id=SI08-P1-CHEM-S4-Q07, source={Score-I - Test 08 (09 Feb 2022)}, paper=1, qno=7, type=numerical, subtopic={Kc and Kp}, answer={2}, answer_src=key, difficulty=6, diff_parts={K1 S3 I1 C1 F1}, diff_note={concepts: first order decomposition  sulphuryl chloride with pressure change | steps: SO2Cl2 gives SO2 + Cl2; total pressure vs time gives extent; first order rate constant from pressure data; partial pressures at time t; compute the asked quantity}]
0.01 moles of sulphuryl chloride $\mathrm{SO_2Cl_2(g)}$ is taken in a sealed tube \& heated to 400 K. It decomposes following 1$^{\text{st}}$ order kinetics according to reaction:

$\mathrm{SO_2Cl_2(g)\rightarrow SO_2(g)+Cl_2(g)}$

The tube is broken after 4 hours \& gas passed through a 25 ml of an acidified 1 N iodine solution where all SO$_2$ is oxidised to $\mathrm{SO_4^{2-}}$. The resulting solution required 10 ml of 1 M Hypo solution. What is half life in hours.
\end{bankq}
\begin{bankq}[id=SI11-P1-CHEM-S3-Q02, source={Score-I - Test 11 (27 Feb 2022)}, paper=1, qno=2, type=numerical, subtopic={Kc and Kp}, answer={0.25}, answer_src=key, difficulty=7, diff_parts={K2 S2 I2 C1 F1}, diff_note={concepts: ratio, standard Gibbs energies at two temperatures from partial pressure plot | steps: K = P_B/P_A from the plot at each temperature; delta G = -R T ln K; ratio of the two values | I: read equilibrium partial pressures from the plot}, also_in={JEE Advanced (2021 Series) - Test 02 (04 Feb 2021) P1 Q2}]
Consider the reaction $\mathbf A\rightleftharpoons\mathbf B$ at 1000 K. At time $t'$, the temperature of the system was increased to 2000 K and the system was allowed to reach equilibrium. Throughout this experiment the partial pressure of \textbf A was maintained at 1 bar. Given below is the plot of the partial pressure of \textbf B with time. What is the ratio of the standard Gibbs energy of the reaction at 1000 K to that at 2000 K?
\bfig{123.5}{SI11_p1_chem_s3_q02_1.png}
\end{bankq}
\begin{bankq}[id=SI24-P1-CHEM-S1-Q01, source={Score-I - Test 24 (11 Dec 2022)}, paper=1, qno=1, type=single, subtopic={Le Chatelier}, answer={A}, answer_src=key, difficulty=6, diff_parts={K2 S2 I2 C1 F0}, diff_note={concepts: enthalpy signs from plots, equilibrium constant against temperature | steps: van t Hoff equation slope of ln K versus 1/T gives minus delta H over R; read slope sign for each plot; compare magnitudes; select | I: interpret graph slopes through the van t Hoff relation}]
Consider following four reactions and given plot showing variation of equilibrium constant ($K_{eq}$) with temperature :

$\mathrm{A+B}\rightleftharpoons\mathrm{C}$\quad $\Delta H_1^\circ$

$\mathrm{D+E}\rightleftharpoons\mathrm{F}$\quad $\Delta H_2^\circ$

$\mathrm{G}\rightleftharpoons\mathrm{H}$\quad $\Delta H_3^\circ$

$\mathrm{I}\rightleftharpoons\mathrm{J}$\quad $\Delta H_4^\circ$
\bfig{108.9}{SI24_p1_chem_s1_q01_1.png}
The correct graph showing the variation of \% yield of product, C, F, H and J with temperature is ($\Delta H^\circ$ is independent of temperature, $\ln K$ is the natural log of equilibrium constant)
\begin{qoptions}
\opt{A}{\bfigi{79.1}{SI24_p1_chem_s1_q01_A.png}}
\opt{B}{\bfigi{77.9}{SI24_p1_chem_s1_q01_B.png}}
\opt{C}{\bfigi{78.2}{SI24_p1_chem_s1_q01_C.png}}
\opt{D}{\bfigi{79.9}{SI24_p1_chem_s1_q01_D.png}}
\end{qoptions}
\end{bankq}
\begin{bankq}[id=SI26-P1-CHEM-S1-Q05, source={Score-I - Test 26 (18 Dec 2022)}, paper=1, qno=5, type=multi, subtopic={Le Chatelier}, answer={A,C,D}, answer_src=key, difficulty=6, diff_parts={K1 S3 I1 C1 F1}, diff_note={concepts: shift of equilibrium for dissolution of ammonium chloride | steps: endothermic dissolution; lowering temperature shifts left; adding chloride common ion shifts left; adding water shifts right; select changes not shifting right}]
The following equilibrium exists in saturated solution of NH$_4$Cl

$\mathrm{NH_4Cl_{(s)}}\rightleftharpoons\mathrm{NH^+_{4(aq)}}+\mathrm{Cl^-_{(aq)}}$

$\Delta H_{25^\circ C}=3.5$ kcal mol$^{-1}$

A change that will not shift the equilibrium to the right is
\begin{qoptions}
\opt{A}{decrease in temperature}
\opt{B}{increase in temperature}
\opt{C}{addition of NH$_4$Cl crystals to the reaction mixture.}
\opt{D}{addition of NH$_4$OH solution to the reaction mixture.}
\end{qoptions}
\end{bankq}
\begin{bankpassage}[id=SI26-P1-CHEM-S3-G1-2, source={Score-I - Test 26 (18 Dec 2022)}, questions={1-2}]
Following equilibria are established on mixing two gases A$_2$ and C.

(i) $3\mathrm{A_{2(g)}}\rightleftharpoons\mathrm{A_{6(g)}}$, \ $K_p=1.6$ atm$^{-2}$

(ii) $\mathrm{A_{2(g)}}+\mathrm{C_{(g)}}\rightleftharpoons\mathrm{A_2C_{(g)}}$

Given that the total pressure is 1.4 atm and partial pressure of A$_6$ is 0.2 atm,

If A$_2$ and C are mixed in 2 : 1 molar ratio, then
\end{bankpassage}
\begin{bankq}[id=SI26-P1-CHEM-S3-Q01, source={Score-I - Test 26 (18 Dec 2022)}, paper=1, qno=1, type=numerical, subtopic={Kc and Kp}, answer={0.40}, answer_src=key, difficulty=5, diff_parts={K1 S1 I2 C1 F1}, diff_note={concepts: Equilibria for A2 A6 and C with reaction coupling | steps: equilibrium partial pressure of C in two coupled equilibria | write Kp for each equilibrium; mass balance on the species; solve the pair of equations for partial pressures | I: two coupled equilibria share a species}, group=SI26-P1-CHEM-S3-G1-2]
Calculate equilibrium partial pressure of C.
\end{bankq}
\begin{bankq}[id=SI26-P1-CHEM-S3-Q02, source={Score-I - Test 26 (18 Dec 2022)}, paper=1, qno=2, type=numerical, subtopic={Kc and Kp}, answer={1.50}, answer_src=key, difficulty=3, diff_parts={K1 S1 I1 C0 F1}, diff_note={concepts: equilibrium constant of the second reaction in atm inverse | steps: Kp(ii) expressed in the partial pressures obtained in the first part; substitute and evaluate | I: depends on equilibrium composition from the previous part}, group=SI26-P1-CHEM-S3-G1-2]
The value of $K_p$ for the reaction (ii) (in atm$^{-1}$) is.
\end{bankq}
\begin{bankq}[id=SI31-P2-CHEM-S1-Q04, source={Score-I - Test 31 (29 Dec 2022)}, paper=2, qno=4, type=single, subtopic={Kc and Kp}, answer={A}, answer_src=key, difficulty=4, diff_parts={K1 S2 I1 C1 F0}, diff_note={concepts: dehydration copper sulphate pentahydrate at different humidity | steps: Kp = p(H2O)^2 with water vapour pressure; relative humidity gives pressure fraction of saturated vapour; compare with Kp at two cities; conclude}]
Consider the following equilibrium at $25^\circ$C

$\mathrm{CuSO_4{\cdot}5H_2O_{(s)}\rightleftharpoons CuSO_4{\cdot}3H_2O_{(s)}+2H_2O_{(g)}}$

If this reaction takes place at two cities A and B having relative humidity 60\% and 80\% respectively. The vapour pressure of water at $25^\circ$C is 20 mm Hg and $K_P$ for the above reaction is $4\times10^{-2}$ (atm)$^2$ at $25^\circ$C then CuSO$_4\cdot$5H$_2$O acts as efflorescent in which city?
\begin{qoptions}
\opt{A}{At A only}
\opt{B}{At B only}
\opt{C}{At both A and B}
\opt{D}{Neither at A nor at B}
\end{qoptions}
\end{bankq}

\section*{Score-II papers}
\begin{bankq}[id=SII01-P2-CHEM-S2-Q09, source={Score-II - Test 01 (21 Mar 2021)}, paper=2, qno=9, type=single, subtopic={Le Chatelier}, answer={B}, answer_src=key, difficulty=3, diff_parts={K1 S1 I1 C1 F0}, diff_note={concepts: match items from physical chemistry lists | steps: compute each value; match}]
Answer the following by appropriately matching the lists based on the information given in the list.

\begin{center}
\begin{tabular}{@{}p{0.40\linewidth}p{0.52\linewidth}@{}}
\textbf{List-I} & \textbf{List-II} \\
(P) NH$_2$COONH$_4$(s) $\rightleftharpoons$ 2NH$_3$(g) $+$ CO$_2$(g); $\Delta H<0$ & (1) High pressure and low temperature will favour formation of products. \\
(Q) CaCO$_3$(s) $\rightleftharpoons$ CaO(s) $+$ CO$_2$(g); $\Delta H>0$ & (2) Increases in the amount of any of the product will definitely favour backward reaction. \\
(R) N$_2$(g) $+$ O$_2$(g) $\rightleftharpoons$ 2NO(g); $\Delta H>0$ & (3) Increase in temperature and doubling the volume of reaction vessel will cause the equilibrium to shift in similar direction. \\
(S) 2SO$_2$(g) $+$ O$_2$(g) $\rightleftharpoons$ 2SO$_3$(g); $\Delta H<0$ & (4) A decrease in temperature will cause decrease in the value of equilibrium constant. \\
 & (5) Increase in concentration of reactant the equilibrium is achieved will make the equilibrium shift towards product side \\
\end{tabular}
\end{center}

Correctly match among the following is:
\begin{qoptions}
\opt{A}{P$\to$3; Q$\to$3,5; R$\to$1,2,5; S$\to$3,4}
\opt{B}{P$\to$2; Q$\to$3,4; R$\to$2,4,5; S$\to$1,2,3,5}
\opt{C}{P$\to$2; Q$\to$4,5; R$\to$1,2,4; S$\to$1,3,5}
\opt{D}{P$\to$2; Q$\to$3,4,5; R$\to$1,4,5; S$\to$2,5}
\end{qoptions}
\end{bankq}
\begin{bankq}[id=SII02-P1-CHEM-S3-Q01, source={Score-II - Test 02 (28 Mar 2021)}, paper=1, qno=1, type=numerical, subtopic={Kc and Kp}, answer={51.20}, answer_src=key, difficulty=6, diff_parts={K1 S3 I1 C1 F1}, diff_note={concepts: equilibrium constant from Gibbs energy at another temperature | steps: delta G0 of the reaction at 300 K; reaction quotient with the given pressures; delta G; K at 300 K; use delta H for another temperature; asked value}]
N$_2$(g) $+$ 3H$_2$(g) $\rightleftharpoons$ 2NH$_3$(g), $\Delta G^0_{f(\text{NH}_3)}=-15.53$ kJ/mole at 300 K. If after 1 hour of mixing reactants, partial pressure of N$_2=50$ Bar, H$_2=2$ Bar, NH$_3=200$ Bar \& $\Delta H_{\text{reaction}}$ is $-4.2$ kJ, find $\Delta S$ of reaction (in JK$^{-1}$). ($\ln10=2.3$, $R=\dfrac{25}{3}$ JK$^{-1}$mol$^{-1}$).
\end{bankq}
\begin{bankq}[id=SII03-P1-CHEM-S1-Q03, source={Score-II - Test 03 (05 May 2021)}, paper=1, qno=3, type=single, subtopic={Kc and Kp}, answer={B}, answer_src=key, difficulty=2, diff_parts={K1 S1 I0 C1 F0}, diff_note={concepts: reactions with Kp greater than Kc | steps: Kp = Kc (RT)^(delta n); delta n positive}]
Select the reactions in which $K_p>K_c$.\\
I. $2\text{NH}_3\underset{}{\overset{T=400K}{\rightleftharpoons}}\text{N}_2(g)+3\text{H}_2(g)$\\
II. $2\text{A}(g)\overset{T=10K}{\rightleftharpoons}2\text{B}(g)+\text{C}(g)$\\
III. $\text{SOCl}_2(g)\overset{450K}{\rightleftharpoons}\text{SO}_2(g)+\text{Cl}_2(g)$\\
IV. $\text{NH}_4\text{NO}_2(s)\overset{370K}{\rightleftharpoons}\text{N}_2(g)+2\text{H}_2\text{O}(\ell)$
\begin{qoptions}
\opt{A}{All}
\opt{B}{I, III, IV}
\opt{C}{I, II, III}
\opt{D}{I \& III}
\end{qoptions}
\end{bankq}
\begin{bankq}[id=SII04-P2-CHEM-S1-Q02, source={Score-II - Test 04 (07 May 2021)}, paper=2, qno=2, type=multi, subtopic={Le Chatelier}, answer={A,B,C}, answer_src=key, difficulty=2, diff_parts={K1 S1 I0 C0 F1}, diff_note={concepts: reactions favouring the forward direction | steps: compare acid base strengths or stability of the products; mark favoured reactions}]
Among the given reaction which of the following favours forward reaction.

\begin{center}
\bfigi{221.5}{SII04_p2chem_q2_reactions.png}
\end{center}
\end{bankq}
\begin{bankq}[id=SII04-P2-CHEM-S1-Q06, source={Score-II - Test 04 (07 May 2021)}, paper=2, qno=6, type=multi, subtopic={Kc and Kp}, answer={A,B,D}, answer_src=key, difficulty=6, diff_parts={K1 S3 I1 C1 F1}, diff_note={concepts: calcium carbonate decomposition equilibrium with added carbon | steps: Kp = P(CO2) = 4 atm; carbon reacts with CO2 to form CO; second equilibrium; pressures; check options}]
CaCO$_3$ decomposed according to the equation $\text{CaCO}_3(s)\rightleftharpoons\text{CaO}(s)+\text{CO}_2(g)$. The equilibrium constant $K_p$ for the reaction is 4 atm. When solid carbon introduced in the vessel a new equilibrium $[\text{CO}_2(g)+\text{C}(s)\rightleftharpoons2\text{CO}(g)]$ is established. The equilibrium constant of this is $2\times10^{-2}$ atm. Then

\begin{qoptions}
\opt{A}{$p_{\text{CO}_2}=4$ atm}
\opt{B}{$p_{\text{CO}}=0.28$ atm}
\opt{C}{At same temperature, if small amount of CaO(s) is added then partial pressure of CO$_2$(g) will decrease.}
\opt{D}{At same temperature, if small amount of CaO(s) is added then partial pressure of CO(g) will remain same.}
\end{qoptions}
\end{bankq}
\begin{bankq}[id=SII09-P1-CHEM-S2-Q10, source={Score-II - Test 09 (12 Feb 2023)}, paper=1, qno=10, type=single, subtopic={Kc and Kp}, answer={B}, answer_src=key, difficulty=3, diff_parts={K1 S2 I0 C1 F0}, diff_note={concepts: stepwise and overall formation constants of a complex | steps: overall constant = product of stepwise constants; concentration of each species at a given ligand concentration; relations among K1 K2 K3 K4}]
Stepwise and overall formation constants ($K_f$) for $[Cu(NH_3)_4]^{2+}$ mentioned as follows:
\resizebox{\linewidth}{!}{$[Cu(H_2O)_4]^{2+}(aq.)+NH_3(aq.)\rightleftharpoons[Cu(NH_3)(H_2O)_3]^{2+}(aq.)+H_2O:K_1$}\par
\resizebox{\linewidth}{!}{$[Cu(NH_3)(H_2O)_3]^{2+}(aq.)+NH_3(aq.)\rightleftharpoons[Cu(NH_3)_2(H_2O)_2]^{2+}(aq.)+H_2O:K_2$}\par
\resizebox{\linewidth}{!}{$[Cu(NH_3)_2(H_2O)_2]^{2+}(aq.)+NH_3(aq.)\rightleftharpoons[Cu(NH_3)_3(H_2O)]^{2+}(aq.)+H_2O:K_3$}\par
\resizebox{\linewidth}{!}{$[Cu(NH_3)_3(H_2O)]^{2+}(aq.)+NH_3(aq.)\rightleftharpoons[Cu(NH_3)_4]^{2+}(aq.)+H_2O:K_4$}\par
\resizebox{\linewidth}{!}{$[Cu(H_2O)_4]^{2+}(aq.)+4NH_3(aq.)\rightleftharpoons[Cu(NH_3)_4]^{2+}(aq.)+4H_2O:K_5$}\par
List-II has the accurate values of the logarithm of the equilibrium constants as described in the above chemical equations. Match the $\log K_f$ in the List-I with their appropriate values given in the List-II.
\begin{center}
\begin{tabular}{|c|c|c|c|}
\hline
\multicolumn{2}{|c|}{List-I} & \multicolumn{2}{c|}{List-II} \\
\hline
(P) & $\log K_1$ & (1) & 2.0 \\
\hline
(Q) & $\log K_2$ & (2) & 2.7 \\
\hline
(R) & $\log K_3$ & (3) & 3.2 \\
\hline
(S) & $\log K_5$ & (4) & 4.0 \\
\hline
 & & (5) & 11.9 \\
\hline
\end{tabular}
\end{center}
\begin{qoptions}
\opt{A}{P$\to$5;Q$\to$4;R$\to$3;S$\to$2}
\opt{B}{P$\to$4;Q$\to$3;R$\to$2;S$\to$5}
\opt{C}{P$\to$2;Q$\to$4;R$\to$3;S$\to$5}
\opt{D}{P$\to$4;Q$\to$3;R$\to$2;S$\to$1}
\end{qoptions}
\end{bankq}
\begin{bankq}[id=SII09-P1-CHEM-S3-Q06, source={Score-II - Test 09 (12 Feb 2023)}, paper=1, qno=6, type=numerical, subtopic={Le Chatelier}, answer={6}, answer_src=key, difficulty=3, diff_parts={K1 S2 I0 C0 F1}, diff_note={concepts: response of an equilibrium to a sudden increase of volume | steps: reaction quotient after doubling the volume; shift toward more moles of gas; compare Kp Kc concentrations total pressure and moles of each species at the new equilibrium}]
The reaction $PCl_5(g)\rightleftharpoons PCl_3(g)+Cl_2(g)$ is at equilibrium. Now suddenly the volume is doubled. The system is now allowed to attain a new equilibrium state. How many of the following are lesser at the new equilibrium as compared to the old equilibrium? Assume the entire process to be isothermal.\\
(a) Moles of $PCl_5$\\
(b) Moles of $PCl_3$\\
(c) Moles of $Cl_2$\\
(d) Concentration of $PCl_5$\\
(e) Concentration of $PCl_3$\\
(f) Concentration of $Cl_2$\\
(g) Rate of the forward reaction\\
(h) Rate of the backward reaction\\
(i) Equilibrium constant
\end{bankq}
\begin{bankq}[id=SII12-P1-CHEM-S4-Q05, source={Score-II - Test 12 (04 Feb 2024)}, paper=1, qno=5, type=numerical, subtopic={Kc and Kp}, answer={4}, answer_src=key, difficulty=5, diff_parts={K1 S2 I1 C1 F1}, diff_note={concepts: equilibrium extent at two temperatures and the van't Hoff relation | steps: read the moles of product at equilibrium for each temperature; Kc for each; relation between K1 and K2 with T1 = 2T2 and delta H; solve for the required quantity}]
In an one-litre flask, 6 moles of A undergoes the reaction $A(g)\rightleftharpoons P(g)$. The progress of product formation at two temperatures (in Kelvin), $T_1$ and $T_2$, is shown in the figure:
\bfig{153.0}{SII12_p1_chem_s4_q05_1.png}
If $T_1=2T_2$ and $(\Delta G_2^\ominus-\Delta G_1^\ominus)=RT_2\ln x$, then the value of x is \_\_\_\_\_\_. [$\Delta G_1^\ominus$ and $\Delta G_2^\ominus$ are standard Gibb's free energy change for the reaction at temperatures $T_1$ and $T_2$, respectively.]
\end{bankq}
\begin{bankq}[id=SII13-P1-CHEM-S1-Q01, source={Score-II - Test 13 (06 Feb 2024)}, paper=1, qno=1, type=single, subtopic={Kc and Kp}, answer={A}, answer_src=key, difficulty=3, diff_parts={K1 S2 I0 C1 F0}, diff_note={concepts: density of a dissociated gas mixture | steps: moles at equilibrium 1.2 with average molar mass 92/1.2 = 76.7; density = P M/RT; evaluate}]
Gaseous $N_2O_4$ dissociates into gaseous $NO_2$ according to the reaction $N_2O_4(g)\rightleftharpoons2NO_2(g)$. At 300 K and 1 atm pressure, the degree of dissociation of $N_2O_4$ is 0.2. If one mole of $N_2O_4$ gas is contained in a vessel, then the density of the equilibrium mixture is :
\begin{qoptions}
\opt{A}{3.11 g/L}
\opt{B}{6.22 g/L}
\opt{C}{1.56 g/L}
\opt{D}{4.56 g/L}
\end{qoptions}
\end{bankq}
\begin{bankq}[id=SII17-P1-CHEM-S1-Q01, source={Score-II - Test 17 (13 Feb 2024)}, paper=1, qno=1, type=single, subtopic={Le Chatelier}, answer={Bonus}, answer_src=solved, working={No answer key available (source truncated) and the figure/data needed to solve this question was not recoverable; banked as Bonus.}, difficulty=3, diff_parts={K1 S2 I1 C0 F0}, diff_note={concepts: equilibrium shift in SO3 preparation with a stated condition | steps: exothermic reaction; Le Chatelier; effect of temperature pressure and catalyst; which statement holds}]
The preparation of SO$_3$(g) by reaction\\
SO$_2$(g) + $\dfrac12$ O$_2$(g) $\rightleftharpoons$ SO$_3$(g)
\bfig{113.0}{SII17_p1_chem_s1_q01_1.png}
is an exothermic reaction. If the preparation follows the following temperature-pressure relationship for its \% yield, then for temperatures $T_1$, $T_2$ and $T_3$. The correct option is :-
\begin{qoptions}
\opt{A}{$T_3>T_2>T_1$}
\opt{B}{$T_1>T_2>T_3$}
\opt{C}{$T_1=T_2=T_3$}
\opt{D}{Nothing could be predicted about temperature through given information.}
\end{qoptions}
\end{bankq}
\begin{bankq}[id=SII19-P1-CHEM-S1-Q08, source={Score-II - Test 19 (20 Feb 2024)}, paper=1, qno=8, type=single, subtopic={Kc and Kp}, answer={A}, answer_src=key, difficulty=4, diff_parts={K1 S2 I1 C1 F0}, diff_note={concepts: equilibrium constant of a combined reaction | steps: add reactions to get the target; K multiplies for sum; K = K1 K2 /K3 type}]
(1) CO$_{(g)}$ + H$_2$O$_{(g)}$ $\rightleftharpoons$ CO$_{2(g)}$ + H$_{2(g)}$ ; K$_1$\\
(2) CH$_{4(g)}$ + H$_2$O$_{(g)}$ $\rightleftharpoons$ CO$_{(g)}$ + 3H$_{2(g)}$ ; K$_2$\\
(3) CH$_{4(g)}$ + 2H$_2$O$_{(g)}$ $\rightleftharpoons$ CO$_{2(g)}$ + 4H$_{2(g)}$ ; K$_3$\\
If the relation between K$_1$, K$_2$ and K$_3$ is $\mathrm{K_3}=\dfrac{(\mathrm{K_1})^x}{(\mathrm{K_2})^y}$ then calculate $(x-y)$
\begin{qoptions}
\opt{A}{2}
\opt{B}{3}
\opt{C}{1}
\opt{D}{0.5}
\end{qoptions}
\end{bankq}
\begin{bankq}[id=SII22-P1-CHEM-S1-Q01, source={Score-II - Test 22 (14 Mar 2024)}, paper=1, qno=1, type=single, subtopic={Kc and Kp}, answer={D}, answer_src=key, difficulty=1, diff_parts={K1 S1 I0 C0 F0}, diff_note={concepts: dependence of Kp on pressure and extent | steps: Kp depends only on temperature; it does not change with P or x}]
At constant temperature, the equilibrium constant ($K_P$) for the decomposition reaction $N_2O_4\rightleftharpoons2NO_2$ is expressed by $K_P=4x^2P/(1-x^2)$ where P is pressure, x is extent of decomposition. Which of the following statement is true ?
\begin{qoptions}
\opt{A}{$K_P$ increases with increase of P}
\opt{B}{$K_P$ increases with increase of x}
\opt{C}{$K_P$ increases with decrease of x}
\opt{D}{$K_P$ remains constant with change in P or x}
\end{qoptions}
\end{bankq}
\begin{bankq}[id=SII24-P1-CHEM-S1-Q07, source={Score-II - Test 24 (17 Mar 2024)}, paper=1, qno=7, type=single, subtopic={Kc and Kp}, answer={A}, answer_src=key, difficulty=4, diff_parts={K1 S2 I1 C1 F0}, diff_note={concepts: equilibrium of water gas shift | steps: ICE table with Kc; extent; composition}]
Carbon monoxide in water gas reacts with steam according to the following reaction:\\
$\mathrm{CO(g)+H_2O(g)\rightleftharpoons CO_2(g)+H_2(g)}$ ; $K_{eq.}=\dfrac13$\\
If two volume of water gas (containing 1 : 1 of CO and H$_2$) are mixed with five volumes of steam, then the volume ratio of the four gases (CO : H$_2$O : CO$_2$ : H$_2$) present at equilibrium is :
\begin{qoptions}
\opt{A}{1 : 9 : 1 : 3}
\opt{B}{1 : 15 : 1 : 5}
\opt{C}{1 : 1 : 1 : 1}
\opt{D}{1 : 3 : 1 : 1}
\end{qoptions}
\end{bankq}
\begin{bankq}[id=SII30-P1-CHEM-S1-Q04, source={Score-II - Test 30 (27 Mar 2024)}, paper=1, qno=4, type=single, subtopic={Kc and Kp}, answer={B}, answer_src=key, difficulty=4, diff_parts={K1 S2 I1 C1 F0}, diff_note={concepts: equilibrium constant from concentrations of the Haber reaction | steps: Kc = [NH3]^2 /([N2][H2]^3); substitute; compare with options}]
The following concentrations were observed at 500 K for the formation of NH$_3$ from N$_2$ and H$_2$. At equilibrium : [N$_2$] = $2\times10^{-2}$ M,\\
\relax[H$_2$] = $3\times10^{-2}$ M and [NH$_3$] = $1.5\times10^{-2}$M. Equilibrium constant for the reaction is \underline{\hspace{2em}}.
\begin{qoptions}
\opt{A}{289}
\opt{B}{417}
\opt{C}{328}
\opt{D}{529}
\end{qoptions}
\end{bankq}
\begin{bankq}[id=SII31-P1-CHEM-S2-Q08, source={Score-II - Test 31 (29 Mar 2024)}, paper=1, qno=8, type=single, subtopic={Kc and Kp}, answer={9}, answer_src=key, difficulty=6, diff_parts={K1 S3 I2 C1 F0}, diff_note={concepts: equilibrium moles from partial pressure with a solid component | steps: set extent x; gas moles total; Kp expression; solve x; the asked number | I: solid B does not appear in Kp}]
A(g) + 2B(s) $\rightleftharpoons$ 2C(g).\\
Initially 2 mol A(g), 4 mole of B(s) and 1 mole of a inert gas are present in a closed container. After equilibrium has established total pressure of container becomes 9 atm. If A(g) consume 50\% at equilibrium then calculate Kp (in atm) for above reaction -
\end{bankq}
\begin{bankpassage}[id=SII36-P1-CHEM-S3-G1-2, source={Score-II - Test 36 (07 Feb 2025)}, questions={1-2}]
Equilibrium of complex formation is of much interest for chemists. Complex formation finds potential applications in qualitative analysis, medicinal chemistry, metallurgy, electroplating etc. Theoretically the formation constant of the complex can be used to infer the stability of the complex. The overall formation constant for the reaction of CN$^-$ with Co$^{2+}$ at 298 K is $10^{19}$ and the rate constant for the formation of [Co(CN)$_6$]$^{3-}$ is $10^3$ M$^{-5}$s$^{-1}$. The standard reduction potential for the reaction [Co(CN)$_6$]$^{3-}$ + e $\to$ [Co(CN)$_6$]$^{4-}$ is $-0.835$ V. Standard reduction potential of Co$^{3+}|$Co$^{2+}$ couple is 1.82 V. $\left[\dfrac{2.303RT}F=0.059\right]$
\end{bankpassage}
\begin{bankq}[id=SII36-P1-CHEM-S3-Q01, source={Score-II - Test 36 (07 Feb 2025)}, paper=1, qno=1, type=numerical, subtopic={Kc and Kp}, answer={19}, answer_src=solved, working={$K_f=k_f/k_d=10^{19}$, so $x=19$.}, difficulty=2, diff_parts={K1 S1 I0 C0 F1}, diff_note={concepts: equilibrium constant as ratio of rate constants | steps: K_f = k_forward / k_backward = 10^19 so x = 19}, group=SII36-P1-CHEM-S3-G1-2]
Rate constant for the formation of [Co(CN)$_6$]$^{4-}$ is $10^X$ times greater than the rate constant for the dissociation of the same complex. Then x =
\end{bankq}
\begin{bankq}[id=SII36-P1-CHEM-S3-Q02, source={Score-II - Test 36 (07 Feb 2025)}, paper=1, qno=2, type=numerical, subtopic={Kc and Kp}, answer={108.4}, answer_src=solved, working={$|\Delta G^\circ|=2.303RT\log K=2.303\times8.314\times298\times19\approx108.4$ kJ/mol.}, difficulty=3, diff_parts={K1 S1 I0 C1 F1}, diff_note={concepts: standard Gibbs energy from the formation constant | steps: delta G = 2.303 R T log K = 2.303 x 8.314 x 298 x 19 J; convert to kJ}, group=SII36-P1-CHEM-S3-G1-2]
The magnitude of standard Gibbs energy change for formation of complex [Co(CN)$_6^{-4}$] (in kJ/mol) is:
\end{bankq}
\begin{bankq}[id=SII39-P2-CHEM-S2-Q06, source={Score-II - Test 39 (02 Mar 2025)}, paper=2, qno=6, type=multi, subtopic={Kc and Kp}, answer={A,C,D}, answer_src=key, difficulty=6, diff_parts={K1 S3 I1 C1 F1}, diff_note={concepts: shifting the equilibrium of ethene hydrogenation | steps: Le Chatelier: exothermic forward so heating shifts left; lower pressure shifts to more moles (left); remove H2 shifts left; add inert gas at constant pressure; check four}]
For the gas phase reaction :\\
C$_2$H$_4$ + H$_2$ $\rightleftharpoons$ C$_2$H$_6$ $(\Delta\text{H}=-32.7\,\text{Kcal})$\\
carried out in a vessel, the equilibrium concentration of C$_2$H$_4$ can be increased by-
\begin{qoptions}
\opt{A}{increasing the temperature}
\opt{B}{decreasing the pressure}
\opt{C}{removing some H$_2$ at constant volume}
\opt{D}{adding some C$_2$H$_6$ at constant volume}
\end{qoptions}
\end{bankq}
\begin{bankq}[id=SII39-P2-CHEM-S3-Q01, source={Score-II - Test 39 (02 Mar 2025)}, paper=2, qno=1, type=numerical, subtopic={Kc and Kp}, answer={6}, answer_src=key, difficulty=8, diff_parts={K2 S3 I1 C2 F1}, diff_note={concepts: standard Gibbs energy at a higher temperature from the degree, dissociation | steps: Kp at 300 K from 50 percent dissociation at 1 bar; delta G at 300 K; delta G0 = delta H0 - T delta S0 needs two data; use given ln relations to find at 1000 K via Van't Hoff with delta H; evaluate P/100 | C: van't Hoff integration for the Gibbs energy at 1000 K}]
At 27$^\circ$C and 1 bar, gaseous N$_2$O$_4$ is 50\% dissociated into NO$_2$ at equilibrium by the reaction :\\
N$_2$O$_4$ $\rightleftharpoons$ 2NO$_2$\\
If magnitude of standard Gibbs free energy change of the reaction at 727$^\circ$C is `P' cal/mol, then find value of $\left(\dfrac P{100}\right)$ ?\\
\relax[Given : ln2 = 0.7, ln3 = 1.1]
\end{bankq}
\begin{bankq}[id=SII39-P2-CHEM-S3-Q03, source={Score-II - Test 39 (02 Mar 2025)}, paper=2, qno=3, type=numerical, subtopic={Kc and Kp}, answer={8}, answer_src=key, difficulty=8, diff_parts={K2 S3 I1 C2 F1}, diff_note={concepts: minimum bromine mass for equilibrium, partial pressures | steps: total pressure 2.25 atm in 164 L at 300 K gives total moles; stoichiometry of BrCl formation from Cl2; Kp from the pressures; liquid bromine must supply BrCl atoms; moles of Br2 consumed; mass | C: partial pressure bookkeeping to find the bromine consumed}]
For the reaction\\
Br$_2$($l$) + Cl$_2$(g) $\rightleftharpoons$ 2BrCl(g) ;\\
In a closed container of volume 164 L initially 10 moles of Cl$_2$ are present at 27$^\circ$C. What minimum mass (in gram) of Br$_2$($l$) must be introduced into this container so that above equilibrium is maintained at total pressure of 2.25 atm. Vapour pressure of Br$_2$($l$) at 27$^\circ$C is 0.25 atm. Assume that volume occupied by liquid is negligible. (R = 0.082 L atm mol$^{-1}$ K$^{-1}$, Atomic mass of Br = 80). Express minimum mass required as $x\times10^2$ g. Determine $x$.
\end{bankq}
\begin{bankq}[id=SII40-P1-CHEM-S4-Q05, source={Score-II - Test 40 (16 Mar 2025)}, paper=1, qno=5, type=numerical, subtopic={Kc and Kp}, answer={16}, answer_src=key, difficulty=5, diff_parts={K1 S2 I1 C1 F1}, diff_note={concepts: equilibrium extent at two temperatures and the van't Hoff relation | steps: read the moles of product at equilibrium for each temperature; Kc for each; relation between K1 and K2 with T1 = 2T2 and delta H; solve for the required quantity}]
In a one-litre flask, 6 moles of A undergoes the reaction A(g) $\rightleftharpoons$ P(g). The progress of product formation at two temperatures (in Kelvin), $T_1$ and $T_2$, is shown in the figure:
\bfig{134.1}{SII40_p1_chem_s4_q05_1.png}
If $T_1=2T_2$ and $\Delta H^0=RT_2\ln x$, then the value of x is \underline{\hspace{1cm}}.
\end{bankq}
\begin{bankq}[id=SII41-P1-CHEM-S1-Q07, source={Score-II - Test 41 (19 Mar 2025)}, paper=1, qno=7, type=single, subtopic={Kc and Kp}, answer={C}, answer_src=key, difficulty=5, diff_parts={K1 S2 I2 C1 F0}, diff_note={concepts: equilibrium pressures with a common product Cl2 | steps: two equilibria share Cl2; Kp expressions with the same partial pressure of chlorine; solve | I: Cl2 is shared so the two Kp expressions couple}]
Two system PCl$_5$(g) $\rightleftharpoons$ PCl$_3$(g) + Cl$_2$(g) and COCl$_2$(g) $\rightleftharpoons$ CO(g) + Cl$_2$(g) are simultaneously in equilibrium in a vessel at constant volume. If some CO is introduced in the vessel, then at new equilibrium, the moles of :
\begin{qoptions}
\opt{A}{PCl$_5$ increase}
\opt{B}{PCl$_5$ remain unchanged}
\opt{C}{PCl$_5$ decrease}
\opt{D}{Cl$_2$ increase}
\end{qoptions}
\end{bankq}
\begin{bankq}[id=SII43-P1-CHEM-S1-Q10, source={Score-II - Test 43 (23 Mar 2025)}, paper=1, qno=10, type=single, subtopic={Kc and Kp}, answer={A}, answer_src=key, difficulty=3, diff_parts={K1 S2 I1 C0 F0}, diff_note={concepts: compound with least equilibrium constant for hydration | steps: bulky and electron donating groups stabilise the carbonyl; ketones lower than aldehydes; rank}]
Which of the following compounds has least equilibrium constant for acid catalyzed hydration :
\begin{qoptions}
\opt{A}{\bfigi{21.2}{SII43_p1_chem_s1_q10_2.png}}
\opt{B}{\bfigi{36.3}{SII43_p1_chem_s1_q10_1.png}}
\opt{C}{\bfigi{44.9}{SII43_p1_chem_s1_q10_3.png}}
\opt{D}{\bfigi{31.8}{SII43_p1_chem_s1_q10_4.png}}
\end{qoptions}
\end{bankq}
\begin{bankq}[id=SII44-P1-CHEM-S4-Q02, source={Score-II - Test 44 (23 Mar 2025)}, paper=1, qno=2, type=numerical, subtopic={Kc and Kp}, answer={1386}, answer_src=key, difficulty=4, diff_parts={K1 S2 I0 C1 F1}, diff_note={concepts: equilibrium constant and standard Gibbs energy | steps: 80 percent dissociation of A to B gives K = 0.8/0.2 = 4 independent of pressure for equal moles; delta G = -RT ln K with R=2 cal and T=500 K; magnitude 1386 cal}]
For a gas phase reaction: A$_{(g)}\rightleftharpoons$B$_{(g)}$, gas A is 80\% dissociated at equilibrium at 10 bar and 227$^\circ$C. Calculate $|\Delta_rG^\circ|$ for the reaction (in cal/mol)?

Given: Rate constant of forward reaction is $0.693\times10^{-3}\ \text{sec}^{-1}$ and $\ell n(2)=0.693$
\end{bankq}
\begin{bankq}[id=SII47-P1-CHEM-S1-Q08, source={Score-II - Test 47 (30 Mar 2025)}, paper=1, qno=8, type=single, subtopic={Kc and Kp}, answer={C}, answer_src=key, difficulty=4, diff_parts={K1 S2 I1 C1 F0}, diff_note={concepts: equilibrium of hydrogen iodide and iodometric titration | steps: degree of dissociation 0.8 for HI with 2 mol each of H2 and I2 gives the equilibrium moles; iodine moles at equilibrium; hypo needed 2 moles per mole of I2; volume at 2 M}]
The degree of dissociation of HI at a particular temperature is 0.8. The volume of 2M hypo solution required to react completely with the iodine present in an equilibrium mixture of reaction, which is started by 2 mole each of H$_2$ and I$_2$ in a vessel of two litre capacity, is
\begin{qoptions}
\opt{A}{0.4 litre}
\opt{B}{0.8 litre}
\opt{C}{1.6 litre}
\opt{D}{3.2 litre}
\end{qoptions}
\end{bankq}

\section*{Score Advanced papers}
\begin{bankq}[id=SA01-P2-CHEM-S3-Q02, source={Score Advanced - Test 01 (22 Sep 2020)}, paper=2, qno=2, type=numerical, subtopic={Kc and Kp}, answer={0.30}, answer_src=key, difficulty=5, diff_parts={K1 S2 I1 C1 F1}, diff_note={concepts: equilibrium of NH4HS with extra ammonia | steps: set up the pressure equation; Kp; solve partial pressure of H2S}]
A definite amount of solid NH$_4$HS is placed in a flask already containing ammonia gas at a fixed temperature and 0.1atm pressure. After the equilibrium is established, total pressure inside the flask is 1.1atm. Value of K$_p$ in (atm)$^2$ is :
\end{bankq}
\begin{bankq}[id=SA04-P1-CHEM-S1-Q02, source={Score Advanced - Test 04 (04 Jul 2021)}, paper=1, qno=2, type=multi, subtopic={Le Chatelier}, answer={A,B,C,D}, answer_src=solved, working={Helium at constant volume leaves all four equilibria undisturbed}, difficulty=4, diff_parts={K1 S1 I1 C1 F1}, diff_note={concepts: adding helium in a rigid vessel | steps: volume is constant so partial pressures do not change}]
The following four reactions are at equilibrium in closed rigid containers. In which of the reaction(s) equilibrium is not disturbed by adding He gas.
\begin{qoptions}
\opt{A}{H$_2$(g) + I$_2$(g) $\rightleftharpoons$ 2HI(g)}
\opt{B}{N$_2$O$_4$(g) $\rightleftharpoons$ 2NO$_2$(g)}
\opt{C}{CO(g) + 2H$_2$(g) $\rightleftharpoons$ CH$_3$OH(g)}
\opt{D}{C(s) + H$_2$O(g) $\rightleftharpoons$ CO(g) + H$_2$(g)}
\end{qoptions}
\end{bankq}
\begin{bankq}[id=SA04-P2-CHEM-S1-Q04, source={Score Advanced - Test 04 (04 Jul 2021)}, paper=2, qno=4, type=single, subtopic={Le Chatelier}, answer={C}, answer_src=solved, working={With dH<0 the slope -dH/R is positive: graph C}, difficulty=5, diff_parts={K2 S1 I2 C1 F0}, diff_note={concepts: van't Hoff graph for an exothermic equilibrium, le chatelier | steps: equilibrium shifts to reactants on heating, so K falls with T; ln K=-dH/RT+dS/R}]
An equilibrium shift towards reactants at higher temperatures. Find the correct graph :
\begin{qoptions}
\opt{A}{\bfigi{55.1}{SA04_p2_chem_s1_q04_1.png}}
\opt{B}{\bfigi{56.0}{SA04_p2_chem_s1_q04_2.png}}
\opt{C}{\bfigi{55.8}{SA04_p2_chem_s1_q04_3.png}}
\opt{D}{\bfigi{56.0}{SA04_p2_chem_s1_q04_4.png}}
\end{qoptions}
\end{bankq}
\begin{bankq}[id=SA04-P2-CHEM-S1-Q08, source={Score Advanced - Test 04 (04 Jul 2021)}, paper=2, qno=8, type=single, subtopic={Kc and Kp}, answer={D}, answer_src=solved, working={Only H2+I2=2HI has dn=0; the others have dn<0}, difficulty=4, diff_parts={K2 S1 I1 C1 F0}, diff_note={concepts: Kc against Kp, kc and kp | steps: Kp=Kc(RT)\^{}dn so Kc<=Kp needs dn>=0}]
The reaction for which K$_\text{C}$ $\leq$ K$_\text{p}$ at 298 K
\begin{qoptions}
\opt{A}{3O$_2$(g) $\rightleftharpoons$ 2O$_3$(g)}
\opt{B}{N$_2$(g) + 3H$_2$(g) $\rightleftharpoons$ 2NH$_3$(g)}
\opt{C}{N$_2$(g) + 2NO$_2$(g) $\rightleftharpoons$ 2NO$_2$(g)}
\opt{D}{H$_2$(g) + I$_2$(g) $\rightleftharpoons$ 2HI(g)}
\end{qoptions}
\end{bankq}
\begin{bankq}[id=SA05-P1-CHEM-S2-Q07, source={Score Advanced - Test 05 (15 Jul 2021)}, paper=1, qno=7, type=numerical, subtopic={Kc and Kp}, answer={-4}, answer_src=key, difficulty=6, diff_parts={K2 S1 I2 C1 F1}, diff_note={concepts: combining equilibrium constants, kc and kp | steps: reverse and halve the first reaction; multiply by the second}]
Consider following gaseous reactions at equilibrium :\par
2AB $\rightleftharpoons$ A$_2$ + B$_2$, \quad K = 2.5 $\times$ 10$^5$\par
AB + $\dfrac{1}{2}$C$_2$ $\rightleftharpoons$ ABC, K = 5 $\times$ 10$^{-3}$\par
If the value of equilibrium constant of following reaction is x $\times$ 10$^y$. Give (x + y) :\par
$\dfrac{1}{2}$A$_2$ + $\dfrac{1}{2}$B$_2$ + $\dfrac{1}{2}$C$_2$ $\rightleftharpoons$ ABC
\end{bankq}
\begin{bankq}[id=SA07-P1-CHEM-S1-Q02, source={Score Advanced - Test 07 (16 Apr 2023)}, paper=1, qno=2, type=multi, subtopic={Kc and Kp}, answer={B}, answer_src=key, difficulty=5, diff_parts={K1 S2 I1 C1 F1}, diff_note={concepts: equilibrium from a concentration graph | steps: stoichiometry from the curves; K expression; effect of changes; choose}]
Attainment of the equilibrium A(g) $\rightleftharpoons$ 2C(g) + B(g) gave the following of graph. Find the correct option(s).\par
\bfig{126.4}{SA07_p1_chem_s1_q02_1.png}
\begin{qoptions}
\opt{A}{Equilibrium const. K$_\text{C}$ for above reaction is 128 M$^2$.}
\opt{B}{At equilibrium \% dissociation of A is 40\%}
\opt{C}{At equilibrium \% dissociation of A is 30\%}
\opt{D}{None of these}
\end{qoptions}
\end{bankq}
\begin{bankq}[id=SA07-P1-CHEM-S3-Q08, source={Score Advanced - Test 07 (16 Apr 2023)}, paper=1, qno=8, type=numerical, subtopic={Kc and Kp}, answer={84.40}, answer_src=key, difficulty=5, diff_parts={K1 S2 I1 C1 F1}, diff_note={concepts: simultaneous dissociation equilibria | steps: partial pressure of each product; Kp for each; total pressure; compute}]
Two solid compounds A and B dissociates into gaseous products at 293 K as\par
(i) A(s) $\rightleftharpoons$ X(g) + H$_2$S(g).\par
(ii) B (s) $\rightleftharpoons$ Y(g) + H$_2$S(g).\par
Under the given condition, pressure over excess solid A when taken alone in a vessel is 68 mm and that over excess solid B when taken alone in a vessel is 50 mm. The total pressure [in mm Hg] if decomposition of both solids A \& B takes place in the same vessel is ......... .
\end{bankq}
\begin{bankq}[id=SA09-P1-CHEM-S1-Q03, source={Score Advanced - Test 09 (29 Apr 2023)}, paper=1, qno=3, type=single, subtopic={Kc and Kp}, answer={C}, answer_src=key, difficulty=5, diff_parts={K1 S3 I1 C1 F0}, diff_note={concepts: density  an equilibrium mixture for ammonia dissociation | steps: M = 17/(1+alpha) average molar mass; density d = PM/RT; relation between alpha and d; equilibrium constant; Kp at 237 C}]
For the equilibrium 2NH$_3$(g) $\rightleftharpoons$ N$_2$(g) + 3H$_2$(g) of reaction starts with pure NH$_3$. Density d(gm/Lit.) of equilibrium mixture varies with degree of dissociation ($\alpha$) of NH$_3$ of 237\textdegree C in such a manner that $\dfrac{1}{\text{d}}$ v/s ``$\alpha$'' is a linear plot shown below, what is the total pressure at equilibrium.\par
\bfig{81.8}{SA09_p1_chem_s1_q03_1.png}
\begin{qoptions}
\opt{A}{5.110 atm}
\opt{B}{3.862 atm}
\opt{C}{2.463 atm}
\opt{D}{1.56 atm}
\end{qoptions}
\end{bankq}
\begin{bankq}[id=SA10-P1-CHEM-S2-Q08, source={Score Advanced - Test 10 (06 May 2023)}, paper=1, qno=8, type=single, subtopic={Le Chatelier}, answer={D}, answer_src=key, difficulty=5, diff_parts={K1 S3 I1 C1 F0}, diff_note={concepts: Le Chatelier on gas phase equilibria | steps: N2O4 on volume increase; H2 + I2 with inert gas at constant P; NO + Br2 on pressure decrease; ICl3 decomposition on pressure increase; match effects | I: inert gas at constant P changes the volume}]
Match the list.\par
{\footnotesize\setlength{\tabcolsep}{1.5pt}\noindent\begin{tabular}{|c|L{3.6cm}|c|L{2.3cm}|}
\hline
 & \multicolumn{1}{c|}{List-I} & & \multicolumn{1}{c|}{List-II}\\ \hline
(P) & N$_2$O$_4$ $\rightleftharpoons$ 2NO$_2$ on increasing volume at constant temperature & (1) & Forward shift.\\ \hline
(Q) & H$_{2(g)}$ + I$_{2(g)}$ $\rightleftharpoons$ 2HI$_{(g)}$ on addition of inert gas at constant P \& T. & (2) & Concentration of reactant decreases.\\ \hline
(R) & 2NO$_{(g)}$ + Br$_{2(g)}$ $\rightleftharpoons$ 2NOBr$_{(g)}$ on decreasing pressure at constant temperature. & (3) & Concentration of product decreases.\\ \hline
(S) & 2ICl$_{3(g)}$ $\rightleftharpoons$ I$_{2(g)}$ + 3Cl$_{2(g)}$ [$\Delta$H $>$ 0] on increasing pressure at constant temperature. & (4) & Equilibrium is not disturbed.\\ \hline
 & & (5) & K$_\text{P}$ $>$ K$_\text{C}$ $\left(\text{T} > \dfrac{1}{\text{R}}\right)$\\ \hline
\end{tabular}}\par
\begin{qoptions}
\opt{A}{P $\longrightarrow$ 1,2,5;Q $\longrightarrow$ 2,3,4;R $\longrightarrow$ 2,3;S $\longrightarrow$ 3,5}
\opt{B}{P $\longrightarrow$ 1,2,5;Q $\longrightarrow$ 1,5;R $\longrightarrow$ 2,3;S $\longrightarrow$ 3,5}
\opt{C}{P $\longrightarrow$ 3;Q $\longrightarrow$ 1,4;R $\longrightarrow$ 1;S $\longrightarrow$ 2}
\opt{D}{P $\longrightarrow$ 1,2,3,5;Q $\longrightarrow$ 2,3,4;R $\longrightarrow$ 2,3;S $\longrightarrow$ 5}
\end{qoptions}
\end{bankq}
\begin{bankq}[id=SA19-P1-CHEM-S3-Q02, source={Score Advanced - Test 19 (04 Apr 2024)}, paper=1, qno=2, type=numerical, subtopic={Kc and Kp}, answer={8}, answer_src=key, difficulty=7, diff_parts={K1 S3 I1 C2 F1}, diff_note={concepts: bromine mass needed to reach equilibrium with Cl2  BrCl | steps: total pressure from moles; Kp = p(BrCl)^2/p(Cl2); partial pressures of Cl2 and BrCl in a 164 L container at 300 K; moles of BrCl formed; bromine consumed; mass | C: translating equilibrium partial pressures into moles with a limiting pure liquid}]
For the reaction: Br$_2$($\ell$) + Cl$_2$(g) $\rightleftharpoons$ 2BrCl(g);$K_p=1$ atm at 300K. In a closed container of volume 164L. initially 10 moles of Cl$_2$(g) are present. If '$m$' gm is the minimum mass of Br$_2$($\ell$) which must be introduced into this container so that above equilibrium is maintained at total pressure of 2.25 atm, then the value of $\left(\dfrac{m}{100}\right)$ is.\par
[Given: vapour pressure of Br$_2$($\ell$) at 27$^\circ$C $=0.25$ atm, volume of liquid Br$_2$, is negligible but sufficient to achieve equilibrium. Temperature of 300K is maintained throughout;\par
$R=0.082$ L-atm / K/ mol; Atomic mass of Br $=80$]\par
Fill your answer as sum of digits (excluding decimal places) till you get the single digit answer.
\end{bankq}
\begin{bankq}[id=SA21-P1-CHEM-S4-Q07, source={Score Advanced - Test 21 (08 Apr 2024)}, paper=1, qno=7, type=numerical, subtopic={Kc and Kp}, answer={5}, answer_src=key, difficulty=7, diff_parts={K1 S3 I1 C2 F1}, diff_note={concepts: simultaneous equilibria sharing nitrogen  hydrogen | steps: two reactions with Kp1 and Kp2; unknown extents; mole balance for N2 H2; partial pressures; solve two equations in two unknowns; required pressure | C: solving two coupled equilibria for the extents}]
In a vessel, two equilibria are simultaneously established at a given temperature as follows:\par
N$_2$(g) + 3H$_2$(g) $\rightleftharpoons$ 2NH$_3$ (g) $[K_{P_1}]$\par
N$_2$(g) + 2H$_2$(g) $\rightleftharpoons$ N$_2$H$_4$ (g) $[K_{P_2}]$\par
Initially the vessel contains N$_2$ and H$_2$ in the molar ratio of 9 : 13. The equilibrium pressure is $7P_0$. The pressure due to NH$_3$ is $P_0$ and due to H$_2$ is $2P_0$. If the value of $(K_{P_1}+K_{P_2})$ is $\dfrac{1}{aP_0^2}$.\par
Find the value of $a$
\end{bankq}
\begin{bankq}[id=SA24-P1-CHEM-S3-Q04, source={Score Advanced - Test 24 (17 Apr 2024)}, paper=1, qno=4, type=numerical, subtopic={Kc and Kp}, answer={6.67}, answer_src=key, difficulty=7, diff_parts={K1 S3 I1 C2 F1}, diff_note={concepts: equilibrium constant  volume or pressure change for PCl5 dissociation | steps: moles 2 2 2 total 6 at pressure 3 atm; partial pressures 1 atm each; Kp = 1; condition for the changed volume or addition; new equilibrium extent; solve | C: re-establishing equilibrium after a perturbation}]
A container of volume `V' litre contains an equilibrium mixture which consists of 2 mol each of PCl$_5$, PCl$_3$ and Cl$_2$. The equilibrium pressure is 3 atm and temperature is `T' kelvin. A certain amount of Cl$_2$ gas is now introduced into the container keeping the temperature and presence constant, until the equilibrium volume becomes `2V' litre. How many moles of Cl$_2$ were added ?
\end{bankq}
\begin{bankq}[id=SA27-P1-CHEM-S3-Q03, source={Score Advanced - Test 27 (24 Apr 2024)}, paper=1, qno=3, type=numerical, subtopic={Kc and Kp}, answer={4}, answer_src=key, difficulty=6, diff_parts={K1 S3 I1 C1 F1}, diff_note={concepts: water gas shift equilibrium  a closed flask | steps: CO + H2O to CO2 + H2; moles 1 each in 2 L; Kc at 700 K; extent x from Kc; equilibrium composition; compute asked}]
A 2.0 L flask, initially containing one mole of each CO and H$_2$O, was sealed and heated to 700 K, where the following equilibrium was established CO (g) + H$_2$O$_{(\mathrm g)}$ $\rightleftharpoons$ CO$_2$ (g) + H$_2$ (g); K$_\mathrm C$ = 9. Now the flask was connected to another flask containing some pure CO$_2$(g) at same temperature and pressure, by means of a narrow tube of negligible volume. When the equilibrium was restored, moles of CO was found to be double of its mole at first equilibrium. Determine volume of CO$_2$(g) flask in litres
\end{bankq}
\begin{bankpassage}[id=SA41-P2-CHEM-S3-G1-2, source={Score Advanced - Test 41}, questions={1-2}]
0.20 moles of CO taken in a 2.463 litre flask is maintained at 750K along with a catalyst so that the following reaction can take place: CO(g) + 2H$_2$(g) $\rightleftharpoons$ CH$_3$OH(g). Hydrogen is introduced until the total pressure of the system is 7.5 atm at equilibrium and 0.1 mole of methanol is formed.\par
(R = 0.0821 L-atm/K-mol).
\end{bankpassage}
\begin{bankq}[id=SA41-P2-CHEM-S3-Q01, source={Score Advanced - Test 41}, paper=2, qno=1, type=numerical, subtopic={Kc and Kp}, answer={0.16}, answer_src=key, difficulty=6, diff_parts={K2 S2 I1 C1 F1}, diff_note={concepts: equilibrium constant Kp for CO + 2H2 forming methanol, pressure data | steps: initial pressures and the final total pressure give extent; partial pressures of CO H2 CH3OH; Kp = P(CH3OH)/(P(CO) P(H2)^2) in atm^-2}, group=SA41-P2-CHEM-S3-G1-2]
For the reaction, K$_\mathrm P$ (in atm$^{-2}$) is
\end{bankq}
\begin{bankq}[id=SA41-P2-CHEM-S3-Q02, source={Score Advanced - Test 41}, paper=2, qno=2, type=numerical, subtopic={Kc and Kp}, answer={12.5}, answer_src=key, difficulty=4, diff_parts={K1 S1 I1 C1 F1}, diff_note={concepts: final pressure without reaction at the same temperature | steps: the same moles of CO and H2 in the same volume exert the sum of the initial partial pressures; apply ideal gas law}, group=SA41-P2-CHEM-S3-G1-2]
What is the final pressure (in atm) if the same amount of CO and H$_2$ as before used, but with no catalyst so that the reaction does not take place?
\end{bankq}

\section*{SRGBT papers}
\begin{bankq}[id=SRGBT05-P1-CHEM-S1-Q19, source={SRGBT - Test 05 (23 Aug 2021)}, paper=1, qno=19, type=multi, subtopic={Kc and Kp}, answer={C}, answer_src=key, difficulty=3, diff_parts={K1 S1 I0 C1 F1}, diff_note={concepts: combining equilibrium constants of added and reversed reactions | steps: target = 2 x (2) - 2 x (1); K = (1.8e4)^2/(0.8)^2 = 5.07e8}]
Given are the information :\par
H$_2$(g) + $\dfrac12$S$_2$(g) = H$_2$S(g) ; K$^\circ_\mathrm P$ = 0.80\par
3H$_2$(g) + SO$_2$(g) = H$_2$S(g) + 2H$_2$O(g) ; K$^\circ_{\mathrm P_2}$ = 1.8$\times$10$^4$\par
The value of K$^\circ_\mathrm P$ for the reaction 4H$_2$(g) + 2SO$_2$(g) = S$_2$(g) + 4H$_2$O(g) is found to be.
\begin{qoptions}
\opt{A}{4.44$\times$10}
\opt{B}{2.25$\times$10$^4$}
\opt{C}{5.07$\times$10$^8$}
\opt{D}{4.26$\times$10$^6$}
\end{qoptions}
\end{bankq}
\begin{bankq}[id=SRGBT05-P1-CHEM-S1-Q20, source={SRGBT - Test 05 (23 Aug 2021)}, paper=1, qno=20, type=multi, subtopic={Kc and Kp}, answer={A}, answer_src=key, difficulty=4, diff_parts={K1 S2 I0 C1 F1}, diff_note={concepts: Kp in terms of the degree of dissociation | steps: moles 1 - alpha and alpha/2 and 3 alpha/2 with total 1 + alpha; Kp expression with P; solve alpha in the given form}]
For the reaction NH$_3$(g) $\rightleftharpoons$ $\dfrac12$N$_2$(g) + $\dfrac32$H$_2$(g), the degree of dissociation in terms K$_\mathrm P$ is--
\begin{qoptions}
\opt{A}{$\alpha=\left(1+\dfrac{3\sqrt3}4\dfrac{\mathrm p}{\mathrm K_\mathrm P}\right)^{-\frac12}$}
\opt{B}{$\alpha=\left(1-\dfrac{3\sqrt3}4\dfrac{\mathrm p}{\mathrm K_\mathrm P}\right)^{-\frac12}$}
\opt{C}{$\alpha=\left(1+\dfrac{3\sqrt3}4\dfrac{\mathrm K_\mathrm P}{\mathrm P}\right)^{\frac12}$}
\opt{D}{$\alpha=\left(1-\dfrac{3\sqrt3}4\dfrac{\mathrm K_\mathrm P}{\mathrm P}\right)^{-\frac12}$}
\end{qoptions}
\end{bankq}
\begin{bankq}[id=SRGBT05-P1-CHEM-S1-Q22, source={SRGBT - Test 05 (23 Aug 2021)}, paper=1, qno=22, type=multi, subtopic={Kc and Kp}, answer={A}, answer_src=key, difficulty=5, diff_parts={K1 S2 I1 C1 F1}, diff_note={concepts: simultaneous equilibria combine to D = C with K = 1 | steps: adding the two reactions gives D = C with K = 1; [C] = [D] = 0.5 nearly; [A] from K1 = [B][C]/[A] with [B] from the second equilibrium gives 1e-6 M}]
1 mole each of A and D is introduced in 1 litre container. Simultaneously the following two equilibrium are established:\par
A $\rightleftharpoons$ B + C; K$_\mathrm C$ = 10$^6$ M\par
B + D $\rightleftharpoons$ A; K$_\mathrm C$ = 10$^{-6}$ M$^{-1}$\par
The equilibrium concentration of A will be
\begin{qoptions}
\opt{A}{10$^{-6}$ M}
\opt{B}{10$^{-3}$ M}
\opt{C}{10$^{-12}$ M}
\opt{D}{10$^{-4}$ M}
\end{qoptions}
\end{bankq}
\begin{bankq}[id=SRGBT05-P1-CHEM-S1-Q24, source={SRGBT - Test 05 (23 Aug 2021)}, paper=1, qno=24, type=multi, subtopic={Kc and Kp}, answer={B}, answer_src=key, difficulty=4, diff_parts={K1 S2 I1 C0 F1}, diff_note={concepts: two heterogeneous equilibria sharing B | steps: Kp1 = PB^2 PC and Kp2 = PB^2 PD; ratio PC/PD = 4; PB = 2(PC + PD) so C false and alpha > alpha'}]
For given two equilibriums attained in a container which are correct if degree of dissociation of A and A$'$ are $\alpha$ and $\alpha'$ respectively (P$'$ = Partial pressure at equilibrium)\par
A(s) $\rightleftharpoons$ 2B(g) + C(g); \quad K$_{\mathrm P_1}$ = 8 $\times$ 10$^{-2}$\par
A$'$(s) $\rightleftharpoons$ 2B(g) + D(g); \quad K$_{\mathrm P_2}$ = 2 $\times$ 10$^{-2}$
\begin{qoptions}
\opt{A}{$\dfrac{\mathrm K_{\mathrm P_2}}{\mathrm K_{\mathrm P_1}}=\left(\dfrac{3\alpha'+2\alpha}{3\alpha+2\alpha'}\right)^3\dfrac\alpha{\alpha'}$}
\opt{B}{$\dfrac{\mathrm P'_\mathrm C}{\mathrm P'_\mathrm D}$ = 4}
\opt{C}{P$'_\mathrm B$ = P$'_\mathrm C$ + P$'_\mathrm D$}
\opt{D}{$\alpha<\alpha'$}
\end{qoptions}
\end{bankq}
\begin{bankq}[id=SRGBT05-P1-CHEM-S1-Q26, source={SRGBT - Test 05 (23 Aug 2021)}, paper=1, qno=26, type=multi, subtopic={Kc and Kp}, answer={A,B,D}, answer_src=key, difficulty=5, diff_parts={K1 S2 I1 C1 F1}, diff_note={concepts: equilibrium ICE table with a large K | steps: reactants in stoichiometric ratio; guess extent x with K = 32^2 x 2/27; solve x and get [C]/[B] and [D]/[A]; dn = 0 so Kp = Kc}]
For a reaction,\par
2A$_{(\mathrm g)}$ + 3B$_{(\mathrm g)}$ $\rightleftharpoons$ 4C$_{(\mathrm g)}$ + D$_{(\mathrm g)}$ \hfill K$^\circ_\mathrm C=\dfrac{32^2\times2}{27}$\par
If 4 moles of A are mixed with 6 moles of B in 2 litre container identify the correct options.
\begin{qoptions}
\opt{A}{At equilibrium, $\dfrac{[\mathrm C]}{[\mathrm B]}=\dfrac83$}
\opt{B}{At equilibrum, $\dfrac{[\mathrm D]}{[\mathrm A]}$ = 1}
\opt{C}{[A] \& [B] will be negligible as compared to [C] \& [D] at equilibrium.}
\opt{D}{K$_\mathrm P$ and K$_\mathrm C$ will be same for the above reaction.}
\end{qoptions}
\end{bankq}
\begin{bankq}[id=SRGBT05-P1-CHEM-S1-Q28, source={SRGBT - Test 05 (23 Aug 2021)}, paper=1, qno=28, type=multi, subtopic={Le Chatelier}, answer={B,C}, answer_src=key, difficulty=4, diff_parts={K2 S2 I0 C0 F1}, diff_note={concepts: Le Chatelier principle for T and inert gas and volume, dG0 versus K | steps: endothermic shifts forward on heating; inert gas at constant volume has no effect; dG0 = 0 only if K = 1; larger volume lowers PB}]
For a gaseous reaction : A$_{(\mathrm g)}$ $\rightleftharpoons$ 3B$_{(\mathrm g)}$ + C$_{(\mathrm g)}$ , $\Delta$H is positive and the reaction attains equilibrium at 1 bar total pressure and 400 K. Identify the incorrect statements regarding the above reaction.
\begin{qoptions}
\opt{A}{On increase of temperature, equilibrium will be shifted in forward reaction.}
\opt{B}{When inert gas is introduced into a rigid container containing above equilibria equilibrium shifts towards left.}
\opt{C}{$\Delta$G$^\circ_{400}$ = 0 for the above reaction.}
\opt{D}{If Volume of vessel containing the above equilibria is increased without change in temperature then partial pressure of B decreases as compared to original equilibrium partial pressure of B.}
\end{qoptions}
\end{bankq}
\begin{bankq}[id=SRGBT05-P1-CHEM-S1-Q31, source={SRGBT - Test 05 (23 Aug 2021)}, paper=1, qno=31, type=multi, subtopic={Kc and Kp}, answer={A,C}, answer_src=key, difficulty=2, diff_parts={K1 S1 I0 C0 F1}, diff_note={concepts: reaction quotient versus K and inert gas effects | steps: Q > K moves backward; inert gas at constant volume has no effect but at constant pressure it does}]
For an equilibrium reaction, which of the following statements is/are correct ?
\begin{qoptions}
\opt{A}{If the reaction quotient of a reaction is greater that K$_\mathrm{eq}$, the reaction has a tendency to move in the backward direction.}
\opt{B}{If the reaction quotient of a reaction is greater than K$_\mathrm{eq}$, the reaction has a tendency to the move in the forward direction.}
\opt{C}{The addition of an inert gas at constant volume does not affect the extent of reaction.}
\opt{D}{The addition of an inert gas at constant pressure does not affect the extent of the reaction in general.}
\end{qoptions}
\end{bankq}
\begin{bankq}[id=SRGBT05-P1-CHEM-S1-Q32, source={SRGBT - Test 05 (23 Aug 2021)}, paper=1, qno=32, type=multi, subtopic={Kc and Kp}, answer={A,C,D}, answer_src=key, difficulty=4, diff_parts={K1 S2 I1 C0 F1}, diff_note={concepts: Le Chatelier shifts propagated through coupled equilibria | steps: Ni removes CO; COCl2 dissociates giving Cl2; extra Cl2 shifts PCl5 back so PCl3 falls and PCl5 rises}]
Following two equilibria are simultaneously established in a container\par
PCl$_5$(g) $\rightleftharpoons$ PCl$_3$(g) + Cl$_2$ ; CO(g) + Cl$_2$(g) $\rightleftharpoons$ COCl$_2$(g)\par
If some Ni(s) is introduced in the container establishing another equilibrium.\par
Ni(s) + 4CO(g) $\rightleftharpoons$ [Ni(CO)$_4$](g)\par
Which among the following statements is/are definitely true about this sytem ?
\begin{qoptions}
\opt{A}{COCl$_2$ concentration decreases as the new equilibrium is established.}
\opt{B}{As Ni is a solid it will not have any impact on the existing two equilibrium reactions.}
\opt{C}{PCl$_3$ concentration decreases to re--attain the equilibrium.}
\opt{D}{PCl$_5$ concentration increases to attain the new equilibrium.}
\end{qoptions}
\end{bankq}
\begin{bankpassage}[id=SRGBT05-P1-CHEM-S1-G35-36, source={SRGBT - Test 05 (23 Aug 2021)}, questions={35-36}]
In a closed rigid vessel, N$_2$ and H$_2$ gases are taken in 9 : 13 mole ratio, by which the following equilibria are estab--lished :\par
N$_2$(g) + 3H$_2$(g) $\rightleftharpoons$ 2NH$_3$(g);K$_{\mathrm P_1}$ = ?\par
N$_2$(g) + 2H$_2$(g) $\rightleftharpoons$ N$_2$H$_4$(g);K$_{\mathrm P_2}$ = ?\par
At equilibrium, the total pressure is 14 atm, the partial pressure of ammonia is 2 atm and the partial pressure of hydrogen is 4 atm.
\end{bankpassage}
\begin{bankq}[id=SRGBT05-P1-CHEM-S1-Q35, source={SRGBT - Test 05 (23 Aug 2021)}, paper=1, qno=35, type=multi, subtopic={Kc and Kp}, answer={B}, answer_src=key, difficulty=3, diff_parts={K1 S1 I0 C1 F1}, diff_note={concepts: Kp from equilibrium partial pressures | steps: partial pressures from the passage data; Kp1 = 1/80 atm^-2}, group=SRGBT05-P1-CHEM-S1-G35-36]
The value of K$_{\mathrm P_1}$ is :
\begin{qoptions}
\opt{A}{$\dfrac1{80}$atm$^2$}
\opt{B}{$\dfrac1{80}$atm$^{-2}$}
\opt{C}{80 atm$^2$}
\opt{D}{$\dfrac1{40}$atm$^{-2}$}
\end{qoptions}
\end{bankq}
\begin{bankq}[id=SRGBT05-P1-CHEM-S1-Q36, source={SRGBT - Test 05 (23 Aug 2021)}, paper=1, qno=36, type=multi, subtopic={Kc and Kp}, answer={B}, answer_src=key, difficulty=2, diff_parts={K1 S1 I0 C0 F1}, diff_note={concepts: equilibrium constant of the reverse reaction | steps: Kp2 from the passage pressures; take the reciprocal 80/3 atm^2}, group=SRGBT05-P1-CHEM-S1-G35-36]
The value of $\dfrac1{\mathrm K_{\mathrm P_2}}$ is :
\begin{qoptions}
\opt{A}{$\dfrac3{80}$atm$^{-2}$}
\opt{B}{$\dfrac{80}3$atm$^2$}
\opt{C}{$\dfrac{40}3$atm$^2$}
\opt{D}{$\dfrac53$atm$^{-2}$}
\end{qoptions}
\end{bankq}
\begin{bankq}[id=SRGBT12-P1-CHEM-S1-Q20, source={SRGBT - Test 12 (10 Oct 2021)}, paper=1, qno=20, type=multi, subtopic={Kc and Kp}, answer={B}, answer_src=key, difficulty=6, diff_parts={K2 S3 I0 C1 F1}, diff_note={concepts: partial pressures after mixing at constant temperature, stoichiometry then dimerisation equilibrium | steps: partial pressures of NO and O2 in 350 mL; NO reacts completely giving NO2 and leftover O2; let x of NO2 dimerise; total pressure 0.365 atm gives x; Kp = P(N2O4)/P(NO2)^2}]
A 250 ml flask containing NO$_\mathrm{(g)}$ at 0.42 atm is connected to a 100 ml flask containing oxygen gas at 0.84 atm by means of a stop cock at 350 K. The gases are mixed by opening the stop cock where the following equilibrium is established.\par
2NO(g) + O$_2$(g)$\longrightarrow$2NO$_2$(g)$\rightleftharpoons$N$_2$O$_4$(g)\par
The first reaction is complete while the second is at equilibrium. If the total pressure is 0.365 atm, K$_\mathrm p$ for dimersation of NO$_2$ is
\begin{qoptions}
\opt{A}{10 atm$^{-1}$}
\opt{B}{0.4 atm$^{-2}$}
\opt{C}{0.16 atm$^{-2}$}
\opt{D}{0.4 atm}
\end{qoptions}
\end{bankq}
\begin{bankq}[id=SRGBT12-P1-CHEM-S1-Q22, source={SRGBT - Test 12 (10 Oct 2021)}, paper=1, qno=22, type=multi, subtopic={Kc and Kp}, answer={A}, answer_src=key, difficulty=5, diff_parts={K1 S2 I1 C1 F1}, diff_note={concepts: heterogeneous equilibria with water vapour (Kp = p(H2O)^n) | steps: vapour pressure of water over each hydrate from the n-th root of Kp; compare with the given humidity; assign each statement T or F}]
Equilibrium constants for the following reactions are given (\textbf{in atm}) for the reactions at 0$^\circ$C. Select the option mentioning correct order of \textbf{True (T)} or \textbf{False (F)} statement.\par
A.6H$_2$O$_{(s)}$ $\rightleftharpoons$ A.2H$_2$O$_{(s)}$ + 4H$_2$O$_{(g)}$ \hfill K$_\mathrm p$ = 1.6 $\times$ 10$^{-11}$\par
B.12H$_2$O$_{(s)}$ $\rightleftharpoons$ B.7H$_2$O$_{(s)}$ + 5H$_2$O$_{(g)}$ \hfill K$_\mathrm p$ = 2.43 $\times$ 10$^{-13}$\par
C.10H$_2$O$_{(s)}$ $\rightleftharpoons$ C$_{(s)}$ + 10H$_2$O$_{(g)}$ \hfill K$_\mathrm p$ = 10$^{-30}$\par
Aqueous tension of H$_2$O at 0$^\circ$C is given as 0.76 torr\par
(I) The most effective drying agent will be C$_{(s)}$ , out of C$_{(s)}$ , B.7H$_2$O$_{(s)}$ \& A.2H$_2$O$_{(s)}$\par
(II) At 0$^\circ$C A.6H$_2$O$_{(s)}$ \& B.12H$_2$O$_{(s)}$ will always be efflorescent.\par
(III) If R.H is less than 100\% in a chamber at 0$^\circ$C then none of the substance can act as deliquescent.
\begin{qoptions}
\opt{A}{TTT}
\opt{B}{FTF}
\opt{C}{TTF}
\opt{D}{TFF}
\end{qoptions}
\end{bankq}
\begin{bankq}[id=SRGBT12-P1-CHEM-S1-Q24, source={SRGBT - Test 12 (10 Oct 2021)}, paper=1, qno=24, type=multi, subtopic={Kc and Kp}, answer={C}, answer_src=key, difficulty=2, diff_parts={K1 S1 I0 C0 F1}, diff_note={concepts: reaction quotient versus K | steps: point P has Q above K = 1.732 so the reaction goes backward}]
For the equilibrium CH$_3$--CH$_2$--CH$_2$--CH$_3$(g) $\rightleftharpoons$ CH$_3$--$\underset{\mathrm H}{\overset{\mathrm{CH}_3}{\mathrm C}}$--CH$_3$(g) equilibrium constant is found to be 1.732 at 298 K. Now if in a vessel at 298 K, a mixture of these two gases be taken as represented by the point P in the figure, predict what will happen :\par
\bfig{56.9}{SRGBT12_p1_chem_s0_q24_1.png}
\begin{qoptions}
\opt{A}{Immediately, above equilibrium will be set up}
\opt{B}{Above reaction will be in the forward direction till it attains equilibrium}
\opt{C}{Above reaction will go in the backward direction till it attains equilibrium}
\opt{D}{Nothing can be said}
\end{qoptions}
\end{bankq}
\begin{bankq}[id=SRGBT12-P1-CHEM-S1-Q26, source={SRGBT - Test 12 (10 Oct 2021)}, paper=1, qno=26, type=multi, subtopic={Kc and Kp}, answer={C,D}, answer_src=key, difficulty=4, diff_parts={K1 S2 I1 C0 F1}, diff_note={concepts: simultaneous equilibria sharing a species | steps: BC is consumed in the second equilibrium so [AB3] > [BC]; B3C2 is formed from BC so [AB3] > [B3C2]; other relations depend on the start}]
If two gases AB$_2$ and B$_2$C are mixed, following equilibria are readily established :\par
AB$_2$(g) + B$_2$C(g)$\rightleftharpoons$AB$_3$(g) + BC(g);\par
BC(g) + B$_2$C(g)$\rightleftharpoons$B$_3$C$_2$(g)\par
If the reaction is started only with AB$_2$ with B$_2$C, then which of the following is necessarily true at equilibrium ?
\begin{qoptions}
\opt{A}{[AB$_3$]$_\mathrm{eq}$ = [BC]$_\mathrm{eq}$}
\opt{B}{[AB$_2$]$_\mathrm{eq}$ $>$ [B$_2$C]$_\mathrm{eq}$}
\opt{C}{[AB$_3$]$_\mathrm{eq}$ $>$ [B$_3$C$_2$]$_\mathrm{eq}$}
\opt{D}{[AB$_3$]$_\mathrm{eq}$ $>$ [BC]$_\mathrm{eq}$}
\end{qoptions}
\end{bankq}
\begin{bankq}[id=SRGBT12-P1-CHEM-S1-Q32, source={SRGBT - Test 12 (10 Oct 2021)}, paper=1, qno=32, type=multi, subtopic={Kc and Kp}, answer={A,B,C}, answer_src=key, difficulty=4, diff_parts={K1 S2 I0 C1 F1}, diff_note={concepts: pressure-temperature scaling with dissociation equilibrium | steps: 20 atm at 400 K becomes 30 atm at 600 K; increase to 50 atm gives alpha = 2/3; partial pressures and Kp = 2.7e3 atm^2}]
Ammonia under a pressure of 20 atm at 127$^\circ$C is heated to 327$^\circ$C in a closed rigid vessel. Under these conditions, ammonia is partially decomposed to N$_2$ and H$_2$ according to the reaction\par
2NH$_3$(g) $\rightleftharpoons$ N$_2$(g) \ + 3H$_2$(g)\par
After decomposition, the pressure inceases to 50 atm. Which of the following statement(s) is/are \textbf{correct} at 327$^\circ$C?
\begin{qoptions}
\opt{A}{The degree of dissociation of NH$_3$ is 2/3}
\opt{B}{The K$_\mathrm p$ of the reaction is 2.7 $\times$ 10$^3$atm$^2$.}
\opt{C}{The partial pressures of both, N$_2$ and NH$_3$ gas at equilibrium is 10 atm.}
\opt{D}{The pressure of NH$_3$ at equilibrium is less than that of N$_2$.}
\end{qoptions}
\end{bankq}
\begin{bankpassage}[id=SRGBT12-P1-CHEM-S1-G35-36, source={SRGBT - Test 12 (10 Oct 2021)}, questions={35-36}]
Following reaction is at equilibrium in basic medium at 300 K.\par
I$_2$ (s) \ $\rightleftharpoons$ I$^-$ (aq.) + IO$_3^-$ (aq.)\par
Equilibrium concentration at 300 K are , [ I$^-$] = 0.1 M, [IO$_3^-$ ] = 0.1 M\par
Given $\Delta$G$^\circ_\mathrm f$ (I$^-$, aq) = $-$ 50 kJ/ mol\par
$\Delta$G$^\circ_\mathrm f$ (IO$_3^-$, aq) = $-$ 123.5 kJ/ mol\par
$\Delta$G$^\circ_\mathrm f$ (H$_2$O, $l$ ) = $-$ 233 kJ/ mol\par
$\Delta$G$^\circ_\mathrm f$ (OH$^-$, aq ) = $-$ 150 kJ/ mol\par
R = $\dfrac{25}3$ J mol$^{-1}$ K$^{-1}$ \ log$_{10}$ e = 2.3
\end{bankpassage}
\begin{bankq}[id=SRGBT12-P1-CHEM-S1-Q35, source={SRGBT - Test 12 (10 Oct 2021)}, paper=1, qno=35, type=multi, subtopic={Kc and Kp}, answer={C}, answer_src=key, difficulty=3, diff_parts={K1 S1 I0 C1 F1}, diff_note={concepts: overall equilibrium constant from the passage data | steps: combine the given steps; K = 1e30}, group=SRGBT12-P1-CHEM-S1-G35-36]
Value of equilibrium constant is
\begin{qoptions}
\opt{A}{10$^{25}$}
\opt{B}{10$^{35}$}
\opt{C}{10$^{30}$}
\opt{D}{10$^{27}$}
\end{qoptions}
\end{bankq}
\begin{bankq}[id=SRGBT12-P1-CHEM-S1-Q36, source={SRGBT - Test 12 (10 Oct 2021)}, paper=1, qno=36, type=multi, subtopic={Kc and Kp}, answer={A}, answer_src=key, difficulty=2, diff_parts={K1 S1 I0 C0 F1}, diff_note={concepts: pH from the equilibrium concentrations | steps: [H+] from the passage equilibrium gives pH 8}, group=SRGBT12-P1-CHEM-S1-G35-36]
pH value at equilibrium is
\begin{qoptions}
\opt{A}{8}
\opt{B}{6}
\opt{C}{9}
\opt{D}{5}
\end{qoptions}
\end{bankq}

\section*{Enthusiast Pre-Score papers}
\begin{bankq}[id=EPS01-P1-CHEM-S2-Q11, source={Enthusiast Pre-Score - Test 01 (04 Nov 2020)}, paper=1, qno=11, type=multi, subtopic={Le Chatelier}, answer={A,D}, answer_src=key, difficulty=6, diff_parts={K1 S3 I1 C1 F1}, diff_note={concepts: equilibrium on raising temperature | steps: exothermic ammonia synthesis; Kp falls; rate of backward; conversion; check}]
For N$_2$(g) + 3H$_2$(g) $\rightleftharpoons$ 2NH$_3$(g)\par
on increasing temperature at 1 atm which of the following is/are correct [Assume enthalpy of reaction is constant]
\begin{qoptions}
\opt{A}{Kp decreases}
\opt{B}{Kp increases}
\opt{C}{Unfavourable $\Delta S_{\mathrm{surr.}}$ increases}
\opt{D}{Favourable $\Delta S_{\mathrm{surr.}}$ decreases}
\end{qoptions}
\end{bankq}
\begin{bankq}[id=EPS01-P1-CHEM-S3-Q02, source={Enthusiast Pre-Score - Test 01 (04 Nov 2020)}, paper=1, qno=2, type=numerical, subtopic={Kc and Kp}, answer={0.165}, answer_src=key, difficulty=5, diff_parts={K1 S2 I1 C1 F1}, diff_note={concepts: equilibrium on change of temperature | steps: equilibrium at 1000 K; Kp at new temperature from van't Hoff; extent; compute}]
Consider the reaction \textbf{A} $\rightleftharpoons$ \textbf{2B} at equilibrium at 1000 K. At time t', the temperature of the system was increased to 3000 K and the system was allowed to reach equilibrium. Throughout this experiment the partial pressure of \textbf{A} was maintained at 1 bar. Given below is the plot of the partial pressure of B with time. What is the ratio of the standard Gibbs energy of the reaction at 1000 K to that at 3000 K?\par
\bfig{154.0}{EPS01_p1_chem_s3_q02_1.png}
\end{bankq}

\section*{JEE Advanced (2021 Series) papers}
\begin{bankq}[id=JA03-P1-CHEM-S3-Q01, source={JEE Advanced (2021 Series) - Test 03 (01 Mar 2021)}, paper=1, qno=1, type=single, subtopic={Kc and Kp}, answer={8.00}, answer_src=key, difficulty=6, diff_parts={K2 S3 I1 C1 F0}, diff_note={concepts: Gibbs energy and equilibrium constant, hydroxide concentration in ionic equilibrium | steps: balance 3 I2 + 6 OH- to 5 I- + IO3- + 3 H2O; delta G0 from formation values; K from delta G0 = -RT ln K with R = 25/3; K expression with I-, IO3- and OH-; solve OH- and pOH; pH = 14 - pOH}]
Calculate the pH at which the following conversion (reaction) will be at equilibrium in basic medium $\mathrm{I_2(s)}\rightleftharpoons\mathrm{I^-(aq)}+\mathrm{IO_3^-(aq)}$ when the equilibrium concentration at 300 K are $(\mathrm{I^-})=0.10$ M and $[\mathrm{IO_3^-}]=0.10$ M.\\
$\{$Given that $\Delta G^0_f(\mathrm{I^-,aq})=-50$ kJ/mole, $\Delta G^0_f(\mathrm{IO_3^-,aq})=-123.5$ kJ/mole, $\Delta G^0_f(\mathrm{H_2O},\ell)=-233$ kJ/mole, $\Delta G^0_f(\mathrm{OH^-,aq})=-150$ kJ/mol, ideal gas constant $=R=\dfrac{25}{3}$ J mole$^{-1}$K$^{-1}$, $\log_e10=2.3\}$
\end{bankq}

\section*{TOAS papers}
\begin{bankq}[id=TOAS06-P1-CHEM-S3-Q01, source={TOAS - Test 06 (17 Apr 2022)}, paper=1, qno=1, type=numerical, subtopic={Kc and Kp}, answer={3.20}, answer_src=key, difficulty=4, diff_parts={K1 S2 I0 C1 F1}, diff_note={concepts: equilibrium and Kp with added gas | steps: initial and equilibrium pressures from the total; Kp; second equilibrium after Cl2 addition}]
Gaseous nitrosyl chloride and nitrogen are taken in a flask sealed and heated to some temperature where the total pressure would have been 1 bar if the following equilibrium had not been established.\\
$\mathrm{2NOCl(g)\overset{K_p}{\rightleftharpoons}2NO(g)+Cl_2(g)}$\\
But the actual pressure was found to be 1.2 bar due to above equilibrium .Now into the above equilibrium mixture, some $\mathrm{Cl_2}$ gas was introduced at constant $V$ \& $T$ so that the total pressure would have been 8.3 bar if no further reaction had occurred but the actual pressure was found to be 8.2 bar. find the value of $k_p$.
\end{bankq}
\begin{bankq}[id=TOAS10-P1-CHEM-S3-Q01, source={TOAS - Test 10 (12 Jun 2022)}, paper=1, qno=1, type=numerical, subtopic={Kc and Kp}, answer={0.02}, answer_src=key, difficulty=2, diff_parts={K1 S1 I0 C0 F1}, diff_note={concepts: equilibrium constant | steps: partial pressures from mole fractions; Kp}]
N$_2$ and H$_2$ gases are mixed in 1 : 3 mole ratio. At equilibrium at 100 atm, the equilibrium mixture is found to contain 20 mole percent NH$_3$. For the reaction :\par
N$_2$(g) + 3H$_2$(g) $\rightleftharpoons$ 2NH$_3$(g)\par
$K_p$ is $\left(\dfrac{x}{216}\right)$ atm$^{-2}$. The value of `x' is.
\end{bankq}
\begin{bankq}[id=TOAS15-P1-CHEM-S3-Q02, source={TOAS - Test 15 (22 Aug 2022)}, paper=1, qno=2, type=numerical, subtopic={Kc and Kp}, answer={729.00}, answer_src=key, difficulty=4, diff_parts={K1 S2 I1 C0 F1}, diff_note={concepts: equilibrium constant from a pressure-time graph | steps: read the equilibrium partial pressures; form the expression with the stoichiometry; evaluate}]
For the dissociation of SO$_3$ into SO$_2$ and O$_2$, at a given temperature, the equilibrium constant $K_p=x$ atm$^2$, then find $x$. The variation of the partial pressure of the three gases with time is as shown :
\bfig{145.4}{TOAS15_p1_chem_s3_q02_1.png}
\end{bankq}
\begin{bankq}[id=TOAS15-P2-CHEM-S2-Q02, source={TOAS - Test 15 (22 Aug 2022)}, paper=2, qno=2, type=numerical, subtopic={Kc and Kp}, answer={0.50}, answer_src=key, difficulty=4, diff_parts={K2 S2 I0 C0 F1}, diff_note={concepts: ideal gas law, degree of dissociation | steps: moles of sulfur vapour from PV = nRT; moles dissociated; fraction}]
2.56gm of sulphur S$_8$(s) is taken which achieves equilibrium with its vapours according to reaction\par
S$_8$(s) $\rightleftharpoons$ 8S(g)\par
If vapours occupies 960 ml at 1atm, 300K, then degree of dissociation of S$_8$(s) will be -\par
[R = 0.08 atm L / mol-K] [Atomic mass : S = 32]
\end{bankq}
\begin{bankpassage}[id=TOAS16-P1-CHEM-S3-G3-4, source={TOAS - Test 16 (04 Sep 2022)}, questions={3-4}]
Consider the following elementary reaction X$_2$(g) $\underset{K_b=4\,\text{sec}^{-1}\text{torr}^{-1}}{\overset{K_f=4\,\text{sec}^{-1}}{\rightleftharpoons}}$ 2X(g) at $2000^\circ$C and at 2 torr equilibrium pressure. The composition of this equilibrium mixture is determined by measuring the rate of effusion through a pin hole. It is observed that the mixture effuses 1.5 times slower than helium under identical conditions.
\end{bankpassage}
\begin{bankq}[id=TOAS16-P1-CHEM-S3-Q03, source={TOAS - Test 16 (04 Sep 2022)}, paper=1, qno=3, type=numerical, subtopic={Kc and Kp}, answer={12}, answer_src=solved, working={$K_p=K_f/K_b=1$ torr $=p_X^2/p_{X_2}$ with $p_X+p_{X_2}=2$ gives $p_X=p_{X_2}=1$, so $\alpha=1/3$. Effusion 1.5 times slower than He: $M_{avg}=4\times2.25=9=M/(1+\alpha)$, $M=12$}, difficulty=6, diff_parts={K2 S2 I1 C1 F1}, diff_note={concepts: equilibrium constant from rate constants, effusion | steps: Kp from kf/kb; equilibrium partial pressures; degree of dissociation; average molar mass from the effusion ratio}, group=TOAS16-P1-CHEM-S3-G3-4]
Calculate molar mass of the gas X$_2$ in gram per mol.
\end{bankq}
\begin{bankq}[id=TOAS16-P1-CHEM-S3-Q04, source={TOAS - Test 16 (04 Sep 2022)}, paper=1, qno=4, type=numerical, subtopic={Kc and Kp}, answer={0.10}, answer_src=key, difficulty=3, diff_parts={K1 S1 I1 C0 F1}, diff_note={concepts: first-order kinetics | steps: fraction reacted at equilibrium; irreversible first-order time}, group=TOAS16-P1-CHEM-S3-G3-4]
If above elementary reaction is assumed to be irreversible then calculate time in sec in which amount of reactant reacted is exactly equal to the amount of reactant reacted to reach equilibrium when reaction was occuring reversibly. $\left(\text{Given : }\ell n\dfrac32=0.4\right)$
\end{bankq}
\begin{bankq}[id=TOAS17-P2-CHEM-S1-Q02, source={TOAS - Test 17 (18 Sep 2022)}, paper=2, qno=2, type=single, subtopic={Kc and Kp}, answer={B}, answer_src=key, difficulty=2, diff_parts={K1 S2 I0 C0 F0}, diff_note={concepts: equilibrium constant in terms of degree of dissociation | steps: moles at equilibrium; partial pressures; expression for Kp}]
For a reaction at equilibrium
\par\noindent\hfill$\mathrm A(g)\rightleftharpoons \mathrm B(g)+\tfrac12\mathrm C(g)$\hfill\null\par
the relation between dissociation constant ($K$), degree of dissociation ($\alpha$) and equilibrium pressure ($p$) is given by :
\begin{qoptions}
\opt{A}{$K=\dfrac{\alpha^{1/2}p^{3/2}}{\left(1+\frac32\alpha\right)^{1/2}(1-\alpha)}$}
\opt{B}{$K=\dfrac{\alpha^{3/2}p^{1/2}}{(2+\alpha)^{1/2}(1-\alpha)}$}
\opt{C}{$K=\dfrac{(\alpha p)^{3/2}}{\left(1+\frac32\alpha\right)^{1/2}(1-\alpha)}$}
\opt{D}{$K=\dfrac{(\alpha p)^{3/2}}{(1+\alpha)(1-\alpha)^{1/2}}$}
\end{qoptions}
\end{bankq}
\begin{bankq}[id=TOAS19-P1-CHEM-S1-Q01, source={TOAS - Test 19 (05 Feb 2023)}, paper=1, qno=1, type=multi, subtopic={Kc and Kp}, answer={A,B,C,D}, answer_src=key, difficulty=3, diff_parts={K1 S2 I0 C0 F1}, diff_note={concepts: chemical equilibrium of phosphorus pentachloride | steps: moles at equilibrium; equilibrium constant; mass; molar mass of the mixture}, also_in={TNAS - Test 01 (05 Feb 2023) P2 Q1}]
In a closed rigid 2 L container the following reaction is occuring :

$\mathrm{PCl_5(g)\rightleftharpoons PCl_3(g)+Cl_2(g)}$

If initially taken 2 moles of PCl$_5$ dissociates to 60\% to establish equilibrium then which of the following is/are correct statement(s) :
\begin{qoptions}
\opt{A}{Mass inside the container remains same.}
\opt{B}{Moles inside the container increases.}
\opt{C}{Molar mass of equilibrium mixture is lesser than molar mass of PCl$_5$.}
\opt{D}{$K_C<1$}
\end{qoptions}
\end{bankq}

\section*{TOPS papers}
\begin{bankq}[id=TOPS05-P2-CHEM-S2-Q05, source={TOPS - Test 05 (14 Jan 2023)}, paper=2, qno=5, type=multi, subtopic={Kc and Kp}, answer={D}, answer_src=key, difficulty=5, diff_parts={K1 S2 I1 C1 F1}, diff_note={concepts: simultaneous equilibria of HI decomposition  another reaction | steps: moles at 127 C; two reactions with extents; pressure or equilibrium constants; solve extents}]
A closed rigid container of volume 1 litre, initially contains 4 mole of HI(g) and 4.25 moles of $\mathrm{N_2(g)}$ at 27$^\circ$C. When the mixture is heated to 127$^\circ$C two reversible reactions takes place till equilibrium is established.

$\mathrm{2HI(g)\rightleftharpoons H_2(g)+I_2(g)}$ \quad $K_{C_1}$

$\mathrm{2N_2(g)+6H_2(g)\rightleftharpoons 4NH_3(g)}$ \quad $K_{C_2}$

At equilibrium the concentration of $\mathrm{I_2}$ and $\mathrm{NH_3}$ are 1.0 M and 0.5 M respectively. Choose the correct option :
\begin{qoptions}
\opt{A}{$K_{C_1}=\dfrac5{32}$}
\opt{B}{$K_{C_2}=\dfrac{16}{125}$}
\opt{C}{$K_{C_1}\times\sqrt{K_{C_2}}=\dfrac1{25}$}
\opt{D}{Conc. of $\mathrm{N_2}$ at equilibrium = 4.0 mol/litre}
\end{qoptions}
\end{bankq}
\begin{bankq}[id=TOPS05-P2-CHEM-S3-Q02, source={TOPS - Test 05 (14 Jan 2023)}, paper=2, qno=2, type=numerical, subtopic={Kc and Kp}, answer={4}, answer_src=key, difficulty=6, diff_parts={K1 S3 I1 C1 F1}, diff_note={concepts: equilibrium composition  SO3 decomposition from titration | steps: two methods give Kc for decomposition and formation; moles at equilibrium; titration of SO2 or acid by 0.2 M reagent; extent; Kc value}]
Equilibrium constant of $\mathrm{SO_3(g)}$ decomposition and formations was calculated by following two methods :
\bfig{112.0}{TOPS05_p2_chem_s3_q02_1.png}
at $t=0$, 10 moles of $\mathrm{SO_3}$ is added 5($\ell$) in container. If at equilibrium, mixture is titrated with 0.2 M acidic $\mathrm{KMnO_4}$ its volume is 4 litre. $K_C$ for the reaction is `$x$' (mol/litre)
\bfig{112.0}{TOPS05_p2_chem_s3_q02_2.png}
Initially 2 : 1 molar ratio $\mathrm{SO_2}$ and $\mathrm{O_2}$ is added respectively. At the equilibrium total pressure of mixture 20 atm. Mole \% of $\mathrm{O_2}$ at equilibrium is 25\%, $K_P$ for the reaction is `$Y$' atm$^{-1}$.

The value of $\dfrac yx$ is ?
\end{bankq}

\section*{TMAS papers}
\begin{bankq}[id=TMAS01-P1-CHEM-S1-Q10, source={TMAS - Test 01 (09 Jul 2023)}, paper=1, qno=10, type=single, subtopic={Le Chatelier}, answer={A}, answer_src=key, difficulty=4, diff_parts={K2 S2 I1 C0 F0}, diff_note={concepts: Le Chatelier, equilibrium | steps: extent of reaction rises with temperature so endothermic; effect of pressure on backward rate; backward rate rises with both | I: rate depends on concentration not equilibrium position}]
For the given reaction :\\
$\mathrm{2A(s)+B(g)\rightleftharpoons C(g)+2D(g)+E(s)}$, in cylinder piston system the extent of reaction of $B$ was found to be 20\% at 300 K and 24\% at 500 K. The rate of backward reaction :
\begin{qoptions}
\opt{A}{Increases with increase in pressure and temperature.}
\opt{B}{Increases with increase in pressure and decrease in temperature.}
\opt{C}{Depends on temperature only and decreases with increases in temperature.}
\opt{D}{Depends on temperature only and increases with increase in temperature.}
\end{qoptions}
\end{bankq}
\begin{bankq}[id=TMAS04-P1-CHEM-S1-Q05, source={TMAS - Test 04 (19 Nov 2023)}, paper=1, qno=5, type=single, subtopic={Kc and Kp}, answer={A}, answer_src=key, difficulty=4, diff_parts={K1 S2 I1 C1 F0}, diff_note={concepts: phase equilibria and Kp relation | steps: boiling point; Kp and Kc; freezing point; solid reaction}]
Select the correct statement(s) from the following option :
\begin{qoptions}
\opt{A}{If water boils at temperature $T^\circ$C at an external pressure of $2$ atm then $[T>100^\circ\mathrm{C}]$}
\opt{B}{For a hypothetical gaseous reaction $\mathrm{A(g)}\rightleftharpoons3\mathrm{B(g)}$, at $10$ Kelvin, $K_P>K_C$ (P in terms of atm and C in terms of mol/litre)}
\opt{C}{Freezing point of water at an external pressure of $2$ atm is greater than its normal freezing point.}
\opt{D}{Position of equilibrium in a reaction, $\mathrm{NH_2COONH_4(s)}\rightleftharpoons2\mathrm{NH_3(g)}+\mathrm{CO_2(g)}$ at a given temperature remains unchanged on addition of more $\mathrm{NH_3(g)}$ at equilibrium.}
\end{qoptions}
\end{bankq}

\section*{TNAS papers}
\begin{bankq}[id=TNAS01-P1-CHEM-S1-Q11, source={TNAS - Test 01 (05 Feb 2023)}, paper=1, qno=11, type=single, subtopic={Kc and Kp}, answer={D}, answer_src=key, difficulty=3, diff_parts={K1 S2 I0 C1 F0}, diff_note={concepts: Kp from partial pressures | steps: mole fractions equal 1/3; partial pressures 4 atm; Kp = 4}]
A reaction at equilibrium involving $2$ moles each $\mathrm{PCl_5}$, $\mathrm{PCl_3}$ and $\mathrm{Cl_2}$ is maintained at $250^\circ$C and a total pressure of $12$ atm. The value of $K_P$ for the reaction

$\dfrac12\mathrm{PCl_3}+\dfrac12\mathrm{Cl_2}\rightleftharpoons\dfrac12\mathrm{PCl_5}$
\begin{qoptions}
\opt{A}{$2$ atm$^{-1/2}$}
\opt{B}{$4$ atm$^{-1/2}$}
\opt{C}{$0.25$ atm$^{-1/2}$}
\opt{D}{$0.5$ atm$^{-1/2}$}
\end{qoptions}
\end{bankq}
\begin{bankq}[id=TNAS01-P1-CHEM-S1-Q13, source={TNAS - Test 01 (05 Feb 2023)}, paper=1, qno=13, type=single, subtopic={Le Chatelier}, answer={A}, answer_src=key, difficulty=2, diff_parts={K1 S2 I0 C0 F0}, diff_note={concepts: Le Chatelier principle | steps: exothermic forward; heating shifts backward; Kc decreases}]
For the reaction : $\mathrm{2SO_2(g)+O_2(g)\rightleftharpoons2SO_3(g)}$; $\Delta H=-$ve. An increase in temperature shows.
\begin{qoptions}
\opt{A}{More dissociation of $\mathrm{SO_3}$ and a decrease in $K_c$}
\opt{B}{Less dissociation of $\mathrm{SO_3}$ and an increase in $K_c$}
\opt{C}{More dissociation of $\mathrm{SO_3}$ and an increase in $K_c$}
\opt{D}{Less dissociation of $\mathrm{SO_3}$ and a decrease in $K_c$}
\end{qoptions}
\end{bankq}
\begin{bankq}[id=TNAS01-P1-CHEM-S2-Q07, source={TNAS - Test 01 (05 Feb 2023)}, paper=1, qno=7, type=numerical, subtopic={Kc and Kp}, answer={9}, answer_src=key, difficulty=6, diff_parts={K2 S2 I1 C1 F1}, diff_note={concepts: simultaneous equilibria, partial pressures | steps: set up the pressure changes for both equilibria; use Kp1 and Kp2; solve | I: common species links the equilibria}]
Pure gas `A' was taken in a closed container at an initial pressure of $5$ atm. The pressure in the container changed due to reactions

$\mathrm{A(g)}\overset{K_{p_1}}{\rightleftharpoons}\mathrm{B(g)}+2\mathrm{C(g)}$

$\mathrm{C(g)}\overset{K_{p_2}}{\rightleftharpoons}\mathrm{B(g)}+\mathrm{D(g)}$

If the total pressure due to mixture of gas A, B, C and D was found to be $10$ atm, in which pressure due to C was $3$ atm. Find $K_{P_1}$ (atm$^2$) for first equilibria.
\end{bankq}
\begin{bankq}[id=TNAS05-P2-CHEM-S1-Q02, source={TNAS - Test 05 (09 Jul 2023)}, paper=2, qno=2, type=multi, subtopic={Le Chatelier}, answer={B,D}, answer_src=key, difficulty=5, diff_parts={K1 S2 I1 C1 F1}, diff_note={concepts: equilibrium and temperature dependence from concentration plots | steps: equilibrium at two temperatures; K comparison; sign of enthalpy; check options | I: K from the final concentrations}]
For the reaction $A(g)\rightleftharpoons P(g)$, the plot of $[A]$ and $[P]$ with time at temperature $T_1$ and $T_2$ are given below. Assume $\Delta H^\circ$ is independent of temperature. If $T_2>T_1$, the correct statement(s) for the reaction $A(g)\rightleftharpoons P(g)$ is/are:
\bfig{89.4}{TNAS05_p2_chem_s1_q02_1.png}
\bfig{88.0}{TNAS05_p2_chem_s1_q02_2.png}
\begin{qoptions}
\opt{A}{$q=\Delta H>0$}
\opt{B}{$W=0$ (at constant pressure)}
\opt{C}{$q=\Delta U>0$}
\opt{D}{$q=\Delta H=\Delta U<0$}
\end{qoptions}
\end{bankq}
\begin{bankq}[id=TNAS07-P1-CHEM-S2-Q06, source={TNAS - Test 07 (20 Aug 2023)}, paper=1, qno=6, type=single, subtopic={Kc and Kp}, answer={9}, answer_src=key, difficulty=8, diff_parts={K2 S3 I2 C2 F0}, diff_note={concepts: degree of dissociation and vapour density, equilibrium constant | steps: stoichiometry from the moles; degree of dissociation from the graph; find n; equilibrium moles; concentrations; Kc}]
In a container of volume 2 dm$^3$, $A(g)$ is allowed to established equilibrium as $A(g)\rightleftharpoons nB(g)$. At equilibrium 5 mole $A$ and 3 mole $B$ are present. Variation of degree of dissociation ($\alpha$) with vapour density is shown by following graph.
\bfig{67.1}{TNAS07_p1_chem_s2_q06_1.png}
where, $D_T$ = Theoretical vapour density, $D_O$ = Observed vapour density.\par
Calculate ($\mathbf{K_c\times10}$) for reaction, $A(g)\rightleftharpoons nB(g)$.
\end{bankq}
\begin{bankq}[id=TNAS07-P2-CHEM-S2-Q09, source={TNAS - Test 07 (20 Aug 2023)}, paper=2, qno=9, type=single, subtopic={Le Chatelier}, answer={D}, answer_src=key, difficulty=6, diff_parts={K1 S3 I2 C1 F0}, diff_note={concepts: equilibrium shifts and Kp Kc relation | steps: each change; shift direction; concentrations; relation of constants; match}]
Match the list.\par
{\footnotesize\setlength{\tabcolsep}{2pt}\begin{tabular}{|c|p{0.38\linewidth}|c|p{0.3\linewidth}|}
\hline
 & \textbf{List-I} & & \textbf{List-II}\\ \hline
(P) & $\mathrm{N_2O_{4(g)}\rightleftharpoons2NO_{2(g)}}$ on increasing volume at constant temperature & (1) & Forward shift.\\ \hline
(Q) & $\mathrm{H_{2(g)}+I_{2(g)}\rightleftharpoons2HI_{(g)}}$ on addition of inert gas at constant P \& T. & (2) & Concentration of reactant decreases.\\ \hline
(R) & $\mathrm{2NO_{(g)}+Br_{2(g)}}$ $\rightleftharpoons$ $\mathrm{2NOBr_{(g)}}$ on decreasing pressure at constant temperature. & (3) & Concentration of product decreases.\\ \hline
(S) & $\mathrm{2ICl_{3(g)}}$ $\rightleftharpoons$ $\mathrm{I_{2(g)}+3Cl_{2(g)}}$ $[\Delta H>0]$ on increasing pressure at constant temperature. & (4) & Equilibrium is not disturbed.\\ \hline
 & & (5) & $K_P>K_C\ \left(T>\dfrac1R\right)$\\ \hline
\end{tabular}}\par
\begin{qoptions}
\opt{A}{$P\to1,2,5;\allowbreak Q\to2,3,4;\allowbreak R\to2,3;\allowbreak S\to3,5$}
\opt{B}{$P\to1,2,5;\allowbreak Q\to1,5;\allowbreak R\to2,3;\allowbreak S\to3,5$}
\opt{C}{$P\to3;\allowbreak Q\to1,4;\allowbreak R\to1;\allowbreak S\to2$}
\opt{D}{$P\to1,2,3,5;\allowbreak Q\to2,3,4;\allowbreak R\to2,3;\allowbreak S\to5$}
\end{qoptions}
\end{bankq}

\section*{Nurture (General) papers}
\begin{bankq}[id=NUR01-P2-CHEM-S2-Q05, source={Nurture (General) - Test 01 (12 Jun 2021)}, paper=2, qno=5, type=multi, subtopic={Kc and Kp}, answer={B}, answer_src=solved, working={Kp=1/(pH2O)\^{}6 gives 5x10\^{}-15 atm; vapour pressure depends only on T; the reaction is exothermic}, difficulty=5, diff_parts={K1 S2 I1 C1 F1}, diff_note={concepts: hydration equilibrium of calcium chloride | steps: equilibrium constant expression with water vapour; vapour pressure; water taken up per formula unit}]
Anhydrous $\mathrm{CaCl_2}$ is often used as desiccant. In the presence of excess $\mathrm{CaCl_2}$ the amount of water taken up is governed by $K_P=6.4\times10^{85}$ for the following reaction at room temperature.\par
$\mathrm{CaCl_2(s)+6H_2O(g)\rightleftharpoons CaCl_2{\cdot}6H_2O(s)}$\par
Which of the following statements is correct about above reaction ?
\begin{qoptions}
\opt{A}{At equilibrium addition of $\mathrm{CaCl_2}$ will shift reaction in forward direction.}
\opt{B}{Vapour pressure of water at equilibrium in a closed vessel that contains $\mathrm{CaCl_2(s)}$ is $5\times10^{-15}$ atm approximately.}
\opt{C}{Vapour pressure of water at equilibrium in a closed vessel becomes half if volume of container is reduced to half at same temperature.}
\opt{D}{If temperature is increased, then vapour pressure of water at equilibrium in closed vessel decreases.}
\end{qoptions}
\end{bankq}
\begin{bankq}[id=NUR04-P1-CHEM-S1-Q02, source={Nurture (General) - Test 04 (11 Feb 2024)}, paper=1, qno=2, type=single, subtopic={Le Chatelier}, answer={B}, answer_src=solved, working={Water is denser than ice, so pressure favours melting}, difficulty=2, diff_parts={K1 S1 I1 C0 F0}, diff_note={concepts: ice water equilibrium under pressure | steps: volume change; Le Chatelier}]
For the equilibrium H$_2$O(s)$\rightleftharpoons$H$_2$O(l)\par
Which of the following statements is true ?
\begin{qoptions}
\opt{A}{The pressure changes do not effect the equilibrium}
\opt{B}{More of ice melts, if pressure on the system is increased}
\opt{C}{More of liquid freezes, if pressure on the system is increased}
\opt{D}{Less of ice melts, if pressure on the system is increased}
\end{qoptions}
\end{bankq}
\end{document}
